\section{후기 (POSTSCRIPT)}\label{uxd6c4uxae30-postscript}

\subsection{서론 (INTRODUCTORY)}\label{uxc11cuxb860-introductory}

\textbf{(p.~238)} 이 책이 처음 출판된 지 32년이 지났다. 그 이후
법리학(jurisprudence)과 철학은 훨씬 더 가까워졌고, 법이론(legal
theory)이라는 주제는 이 나라와 미국 양쪽에서 크게 발전하였다. 학계의
법률가들과 철학자들 사이에서 이 책의 주요 교설들(doctrines)에 대한
비판자들이 그 교설들에 동조한 이들만큼이나 많았더라도, 나는 이 책이
이러한 발전을 자극하는 데 어느 정도 도움을 주었다고 생각하고 싶다.
어쨌든, 내가 원래 이 책을 영국 학부생 독자들을 염두에 두고 썼지만, 이
책은 훨씬 더 넓게 유통되었고 영어권 세계와 이 책의 번역본들이 나온 여러
나라에서 비판적 논평의 방대한 부차 문헌을 낳았다. 이 비판 문헌의 많은
부분은 법학 및 철학 저널의 논문들로 이루어져 있지만, 그 밖에도 이 책의
여러 교설들이 비판의 표적이자 비판자들 자신의 법이론 전개의 출발점으로
취급된 중요한 책들이 여럿 출판되었다.

나는 내 비판자들 가운데 몇몇, 특히 고 론 풀러 교수\footnote{그의
  \emph{The Morality of Law} (1964)에 대한 나의 서평을 보라,
  \emph{Harvard Law Review} 제78권 (1965) 1281면. 이 글은 \emph{Essays
  in Jurisprudence and Philosophy} (1983)에도 재수록됨, 343면. {[}주:
  대괄호로 표시된 각주는 편집자들이 추가한 것임.{]}}와 R. M. 드워킨
교수\footnote{나의 다음 글을 참조하라: “Law in the Perspective of
  Philosophy: 1776--1976”, \emph{New York University Law Review} 제51권
  (1976) 538면; “American Jurisprudence through English Eyes: The
  Nightmare and the Noble Dream”, \emph{Georgia Law Review} 제11권
  (1977) 969면; “Between Utility and Rights”, \emph{Columbia Law
  Review} 제79권 (1979) 828면. 이들 모두 \emph{Essays in Jurisprudence
  and Philosophy}에 재수록되어 있다. 또한, “Legal Duty and
  Obligation”, \emph{Essays on Bentham} (1982)의 제6장, 그리고 R.
  Gavison 편, \emph{Issues in Contemporary Legal Philosophy} (1987) 수록
  “Comment”(35면)를 참조하라.}를 향해 몇 발의 견제 사격을 가한 적은
있지만, 지금까지 그들 중 누구에게도 일반적이고 포괄적인 답변을 하지는
않았다. 나는 일부 비판자들이 \textbf{(p.~239)} 나와 달랐던 만큼이나
서로도 달랐던, 매우 유익한 계속 논쟁을 지켜보고 거기서 배우는 편을
선호했다. 그러나 이 후기에서는 드워킨이 \emph{Taking Rights Seriously}
(1977)와 \emph{A Matter of Principle} (1985)에 수록된 여러 선구적 논문들
및 그의 책 \emph{Law's Empire} (1986)\footnote{이하에서는 각각
  \emph{TRS}, \emph{AMP}, \emph{LE}로 인용함.}에서 제기한 광범위한
비판들 가운데 일부에 답하려고 한다. 이 후기에서 나는 주로 드워킨의
비판들에 초점을 맞춘다. 그는 이 책의 거의 모든 독특한 테제들이
근본적으로 잘못되었다고 논변했을 뿐 아니라, 이 책에 암묵적으로 들어 있는
법이론의 전체 구상(conception)과 법이론이 무엇을 해야 하는지의 전체
구상을 문제 삼았기 때문이다. 이 책의 주요 주제들에 대한 드워킨의
논변들은 여러 해 동안 대체로 일관되었지만, 몇몇 논변들의 실질과 그것들이
표현되는 용어법 양쪽에는 중요한 변화들이 있었다. 그의 초기 논문들에서
두드러졌던 몇몇 비판들은, 명시적으로 철회되지는 않았지만, 그의 후기
작업에는 나타나지 않는다. 그러나 그러한 초기 비판들은 널리 통용되고 매우
영향력이 있으므로, 나는 그의 후기 비판들뿐 아니라 그것들에도 답하는 것이
적절하다고 생각했다. 이 후기의 첫째이자 더 긴 부분은 드워킨의 논변들을
다룬다. 그러나 둘째 부분에서는, 내 테제들 가운데 일부에 대한 나의 전개에
불명료함과 부정확성이 있을 뿐 아니라 어떤 지점들에서는 실제 비일관성과
모순도 있다는 여러 다른 비판자들의 주장들을 검토한다.\footnote{{[}하트는
  본 절에서 언급된 두 개 절 중 두 번째를 완성하지 못함. 편집자 주
  참조.{]}} 여기서 나는, 생각하고 싶지 않을 만큼 많은 경우에 내
비판자들이 옳았음을 인정해야 하며, 이 후기를 기회로 삼아 불명료한 것을
명료화하고 내가 원래 쓴 것 가운데 비일관적이거나 모순적인 곳을
수정하고자 한다.

\subsection{1. 법이론의 본성 (THE NATURE OF LEGAL
THEORY)}\label{uxbc95uxc774uxb860uxc758-uxbcf8uxc131-the-nature-of-legal-theory}

이 책에서 나의 목표는 일반적이면서도 기술적인(descriptive), 법이
무엇인가에 관한 이론을 제공하는 것이었다. 그것이 \emphksans{일반적}이라는
것은, 그것이 어떤 특정 법체계나 법문화에 묶여 있지 않고,
규칙-지배적이고(rule-governed) 그 의미에서 ‘규범적(normative)’인 측면을
가진 복합적 사회·정치 제도로서의 법에 관하여 설명적이고 명료화하는
설명(account)을 제시하려 한다는 의미이다. 이 \textbf{(p.~240)} 제도는
서로 다른 문화와 서로 다른 시대의 많은 변이에도 불구하고 동일한 일반적
형식과 구조를 취해 왔지만, 명료화를 요구하는 많은 오해들과 흐리게 하는
신화들이 그 주변에 모여 있었다. 이 명료화 작업의 출발점은, 이 책 3쪽에서
내가 어느 교육받은 사람에게나 귀속시킨 근대적 국내법체계의 두드러진
특징들에 관한 널리 퍼진 공통 지식이다. 나의 설명(account)은 도덕적으로
중립적이고 정당화적 목적들을 가지지 않는다는 점에서 \emphksans{기술적}이다.
그것은 법에 관한 나의 일반적 설명에 나타나는 형식들과 구조들을 도덕적
근거 또는 다른 근거 위에서 정당화하거나 권장하려 하지 않는다. 다만 나는
이것들에 대한 명료한 이해가 법에 대한 어떤 유익한 도덕적 비판에도 중요한
예비 단계라고 생각한다.

이 기술적 기획을 수행하는 수단으로서, 내 책은 \emphksans{책무부과
규칙들(duty-imposing rules), 권한부여 규칙들(power-conferring rules),
승인 규칙들(rules of recognition), 변경 규칙들(rules of change),
규칙들의 수용(acceptance of rules), 내부 관점과 외부 관점(internal and
external points of view), 내부 진술들과 외부 진술들(internal and
external statements)}, 그리고 \emphksans{법적 유효성(legal validity)} 같은
여러 개념들을 반복해서 사용한다. 이 개념들은 여러 법제도들과 법적
실천들이 해명적으로 분석될 수 있는 항들이 되는 요소들에 주의를
집중시키고, 이러한 제도들과 실천들에 관한 성찰이 촉발한 법의 일반적
본성에 관한 질문들에 해답들이 제시될 수 있게 한다. 여기에는 다음과 같은
질문들이 포함된다. 규칙들이란 무엇인가? 규칙들은 단순한 습관들이나
행위의 반복적 일양성들과 어떻게 다른가? 법규칙들에는 근본적으로 서로
다른 유형들이 있는가? 규칙들은 어떻게 관련될 수 있는가? 규칙들이 하나의
체계를 형성한다는 것은 무엇인가? 법규칙들과 그것들이 가진 권위는
한편으로 위협들과, 다른 한편으로 도덕적 요구사항들과 어떻게
관련되는가?\footnote{H. L. A. Hart, “Comment”, 위의 Gavison 편, 주석
  2, 35면.}

이런 방식으로 기술적이면서 일반적인 것으로 구상된 법이론은, 드워킨이
법이론을, 또는 그가 자주 부르듯 ‘법리학(jurisprudence)’을, 부분적으로
평가적이고 정당화적이며 ‘특정 법문화에 향한(addressed to a particular
legal culture)’ 것으로 보는 구상과 근본적으로 다른 기획이다.\footnote{\emph{LE}
  102면.} 그 법문화는 대개 이론가 자신의 법문화이고, 드워킨의 경우에는
영미법의 법문화이다. 이렇게 구상된 법이론의 중심 과제는 드워킨에 의해
‘해석적(interpretive)’이라고 불리며\footnote{\emph{LE} 제3장.}
부분적으로 평가적이다. 왜냐하면 그것은 한 법체계의 확립된 법(settled
law)과 법적 실천들에 가장 잘 ‘부합(fit)’하거나 그것들과 정합하는(cohere)
동시에 그것들에 대한 최선의 도덕적 정당화를 제공하는 \textbf{(p.~241)}
원칙들을 식별하는 데, 그리하여 법을 ‘그 최선의 빛 아래’ 보이는 데 있기
때문이다.\footnote{\emph{LE} 90면.} 드워킨에게 이렇게 식별된 원칙들은
법에 관한 이론의 일부일 뿐 아니라 법 자체의 암시적 부분들이기도 하다.
그래서 그에게 “법리학은 재판(adjudication)의 일반적 부분이며, 법에 의한
어떤 결정(any decision at law)에나 붙는 침묵의 서곡”이다.\footnote{\emph{LE}
  90면.} 그의 초기 작업에서 그러한 원칙들은 단순히 “법에 관한 가장
건전한 이론(the soundest theory of law)”으로 지칭되었지만,\footnote{\emph{TRS}
  66면.} 그의 최신 작업인 \emph{Law's Empire}에서 그는 이러한 원칙들과
그것들에서 따라 나오는 특정 법명제들(propositions of law)을 ‘해석적
의미’에서의 법으로 특징짓는다. 그러한 해석이론이 해석해야 할 확립된 법적
실천들 또는 법의 전형들(paradigms)은 드워킨에 의해
‘전해석적(preinterpretive)’이라고 불리며,\footnote{\emph{LE} 65--66면.}
이론가는 그러한 전해석적 자료들을 식별하는 데 어려움도 없고 수행해야 할
이론적 과제도 없는 것으로 여겨진다. 그것들은 특정 법체계들의 법률가들의
일반적 합의의 문제로 확립되어 있기 때문이다.\footnote{그러나 드워킨은
  그러한 전해석적 법(preinterpretive law)을 식별하는 일 자체가 해석을
  수반할 수 있다고 경고한다. \emph{LE} 66면 참조.}

나 자신의 법이론 구상과 드워킨의 법이론 구상처럼 그렇게 서로 다른 기획들
사이에 왜 중대한 충돌이 있어야 하는지, 또는 기실 있을 수 있는지는
분명하지 않다. 따라서 \emph{Law's Empire}를 포함한 드워킨 작업의 많은
부분은, 법, 곧 ‘과거의 정치적 결정들’\footnote{\emph{LE} 93면.}이 강제를
정당화하는 방식에 관한 세 가지 서로 다른 설명들의 비교상 이점들을
정교하게 전개하는 데 바쳐져 있으며, 그로부터 그가
‘관행주의(conventionalism)’, ‘법적 실용주의(legal pragmatism)’,
‘통합성으로서의 법(law as integrity)’이라고 부르는 세 가지 서로 다른
법이론 형식들이 나온다.\footnote{\emph{LE} 94면.} 그가 이 세 유형의
이론에 관해 쓰는 모든 것은 평가적·정당화적 법리학에 대한 기여로서 큰
흥미와 중요성을 가진다. 나는, 이 책에서 제시된 것과 같은 실증주의적
법이론이 그러한 해석이론으로 해명적으로 재진술될 수 있다고 그가 주장하는
한을 제외하고, 이러한 해석적 관념들을 그가 전개한 것에 이의를 제기하려는
것이 아니다.\footnote{그러나 일부 비평가들, 예컨대 Michael Moore는
  “\,`The Interpretive Turn in Modern Theory: A Turn for the
  Worse?'\,”, \emph{Stanford Law Review} 제41권 (1989) 871면에서
  947--948면, 드워킨의 의미에서 법실천이 해석적이라는 점에는
  동의하면서도, 법이론은 해석적일 수 없다고 주장한다.} 이 후자의 주장은
\textbf{(p.~242)} 내 견해로는 잘못되었으며, 아래에서 내 이론의 그러한
해석적 판본에 반대하는 이유들을 제시한다.

그러나 드워킨은 그의 책들에서 일반적이고 기술적인 법이론을 잘못 이끌린
것이거나, 기껏해야 단순히 쓸모없는 것으로 배제하는 듯하다. 그는 “유용한
법이론들”은 “역사적으로 발전하는 실천의 특정 단계에 대한 해석”이라고
말하며,\footnote{\emph{LE} 102면. 또한, “우리에게 법의 일반 이론들은
  우리 자신의 사법 실천에 대한 일반 해석들이다”라는 진술도 비교하라.
  \emph{LE} 410면.} 이전에는 “기술(description)과 평가(evaluation)
사이의 평면적 구별(flat distinction)”이 “법이론을 약화시켰다”고
썼다.\footnote{\emph{AMP} 148면. 또한, “\,‘법이론들은 사회적 실천에
  대한 중립적 설명들로 \ldots{} 합당하게 이해될 수 없다’\,”는 진술도
  비교하라. Ronald Dworkin, “A Reply by Ronald Dworkin”, Marshall
  Cohen 편, \emph{Ronald Dworkin and Contemporary Jurisprudence} (1983),
  이하 \emph{RDCJ}로 인용, 247면 중 254면.}

나는 드워킨이 기술적 법이론, 또는 그가 자주 부르듯 ‘법리학’을 거부하는
정확한 이유들을 따라가기 어렵다. 그의 중심 반론은, 법이론은 하나의
법체계 안의 내부자 또는 참여자의 관점인 법에 대한 내부 관점(internal
perspective)을 고려해야 하며, 참여자의 관점이 아니라 외부 관찰자의
관점을 가진 기술 이론으로는 이 내부 관점에 대한 적절한 설명이 제공될 수
없다는 것처럼 보인다.\footnote{{[}\emph{LE} 13--14면 참조.{]}} 그러나 내
책에서 예시된 기술적 법리학의 기획에는, 비참여자 외부 관찰자가
참여자들이 그러한 내부 관점에서 법을 바라보는 방식들을 기술하는 일을
배제하는 것이 사실상 전혀 없다. 그래서 나는 이 책에서 참여자들이 법을
자기 행위에 대한 지침들과 비판 기준들을 제공하는 것으로 수용함으로써
자신의 내부 관점을 드러낸다고 상당히 길게 설명했다. 물론 기술적
법이론가는 그러한 자격으로는 참여자들이 이러한 방식으로 법을 수용하는
것을 스스로 공유하지 않는다. 그러나 그는 그러한 수용을 기술할 수 있고 또
기술해야 하며, 기실 나는 이 책에서 그렇게 하려고 시도하였다. 이 목적을
위해 기술적 법이론가는 내부 관점을 채택한다는 것이 무엇인지
\emphksans{이해}해야 하고, 그러한 제한된 의미에서 그는 내부자의 자리에 자신을
둘 수 있어야 한다는 것은 참이다. 그러나 이것은 법을 수용하거나, 내부자의
내부 관점을 공유하거나 지지(endorse)하거나, 어떤 다른 방식으로든 자신의
기술적 입장을 포기하는 것이 아니다.

드워킨은 기술적 법리학에 대한 자신의 비판에서, 외부 관찰자가 참여자의
내부 관점을 이런 기술적 방식으로 고려할 수 있다는 \textbf{(p.~243)} 이
명백한 가능성을 배제하는 듯하다. 앞서 말했듯이 그는 법리학을
‘재판(adjudication)의 일반적 부분’과 동일시하고, 이것은 법리학 또는
법이론 자체를 그 사법적 참여자들의 내부 관점에서 보인 한 체계의 법의
일부로 취급하는 것이기 때문이다. 그러나 기술적 법이론가는 법에 관한
내부자의 내부 관점을 채택하거나 공유하지 않고도 그것을 이해하고 기술할
수 있다. 설령 (닐 매코믹\footnote{{[}\emph{Legal Reasoning and Legal
  Theory} (1978), 63--64면 및 139--140면 참조.{]}}과 많은 다른
비판자들이 논변했듯이) 행위 지침들과 비판 기준들을 제공하는 것으로 법을
수용하는 데서 드러나는 참여자의 내부 관점이, 법의 요구사항들에 부합할
\emphksans{도덕적} 이유들이 있고 법의 강제 사용에 대한 \emphksans{도덕적} 정당화가
있다는 믿음도 필연적으로 포함한다 하더라도, 이것 역시 도덕적으로
중립적인 기술적 법리학이 기록할 것이지 지지하거나 공유할 것은 아니다.

그러나 드워킨이 ‘해석적’이라고 부르는 부분적으로 평가적인 쟁점들이
법리학과 법이론의 유일하게 적절한 쟁점들이 아니며 일반적이고 기술적인
법리학에도 중요한 자리가 있다는 나의 주장에 대해, 그는 이것이 그렇다는
점을 인정하였다. 그리고 그는 ‘법리학은 재판의 일반적 부분’이라는 자신의
관찰들과 같은 것들이 한정될 필요가 있다고 설명하였다. 그가 이제 말하듯,
이것은 오직 ‘뜻의 문제(question of sense)에 관한 법리학에 대해서만
참’이기 때문이다.\footnote{R. M. Dworkin, “Legal Theory and the Problem
  of Sense”, R. Gavison 편, \emph{Issues in Contemporary Legal
  Philosophy: The Influence of H. L. A. Hart} (1987), 19면.} 이것은
법이론의 유일하게 적절한 형식은 해석적이고 평가적인 것이라는 과도하고,
기실 드워킨 자신이 표현했듯 ‘제국주의적’인 주장으로 보였던 것에 대한
중요하고 환영할 만한 교정이다.

그러나 드워킨이 이제 자신의 겉보기에 제국주의적인 주장의 철회와 결부시킨
다음의 주의를 주는 말들(cautionary words)의 함의는 나에게 여전히 매우
당혹스럽다. “그러나 하트의 것과 같은 일반 이론들이 주로 논의해 온
쟁점들에서 그 {[}뜻의{]} 문제가 얼마나 널리 퍼져 있는지를 강조할 가치가
있다.”\footnote{앞 인용과 동일.} 이 주의의 관련성은 분명하지 않다. 내가
논의한 쟁점들(위 p.~240의 목록 참조)에는 \textbf{(p.~244)} 한편으로
강제적 위협들과, 다른 한편으로 도덕적 요구사항들과 법의 관계 같은
질문들이 포함된다. 드워킨의 주의의 요점은, 그러한 쟁점들을 논의할 때
기술적 법이론가조차 법명제들(propositions of law)의 뜻(sense) 또는
의미(meaning)에 관한 질문들에 직면해야 하며, 이 질문들은 해석적이고
부분적으로 평가적인 법이론에 의해서만 만족스럽게 답될 수 있다는 것처럼
보인다. 이것이 정말 사실이라면, 어떤 주어진 법명제의 뜻을 결정하기 위해
기술적 법이론가조차 “이 명제가 확립된 법에 가장 잘 부합하고 그것을 가장
잘 정당화하는 원칙들에서 따라 나오려면, 이 명제에는 어떤 의미가
배정되어야 하는가?”라는 해석적이고 평가적인 질문을 묻고 답해야 할
것이다. 그러나 내가 언급한 종류의 질문들에 대한 답을 찾는 일반적이고
기술적인 법이론가가 많은 서로 다른 법체계들에서 법명제들의 의미를
결정해야 한다는 것이 참이라 하더라도, 이것이 드워킨의 해석적이고
평가적인 질문을 묻는 일에 의해 결정되어야 \emphksans{한다}는 견해를 받아들일
이유는 없어 보인다. 더구나 일반적이고 기술적인 법이론가가 고려해야 하는
모든 법체계들의 판사들과 법률가들 자신이 사실상 의미의 질문들을 이러한
해석적이고 부분적으로 평가적인 방식으로 확정한다고 하더라도, 이것은
그러한 법명제들의 의미에 관한 자기의 일반적 기술적 결론들을 기초 지을
하나의 사실로서 일반적 기술 이론가가 기록해야 할 일일 것이다. 이러한
결론들이 그렇게 기초 지어졌다는 이유로, 그 결론들 자체가 해석적이고
평가적이어야 하며, 그것들을 제시하면서 이론가가 기술의 과제에서 해석과
평가의 과제로 옮겨 갔다고 상정하는 것은 물론 중대한 오류일 것이다.
기술(description)은 기술되는 것이 평가(evaluation)일 때에도 여전히
기술일 수 있다.

\subsection{2. 법실증주의의 본성 (THE NATURE OF LEGAL
POSITIVISM)}\label{uxbc95uxc2e4uxc99duxc8fcuxc758uxc758-uxbcf8uxc131-the-nature-of-legal-positivism}

\subsubsection{\texorpdfstring{(i) \emphksans{의미론적 이론으로서의 실증주의}
(Positivism as a Semantic
Theory)}{(i) 의미론적 이론으로서의 실증주의 (Positivism as a Semantic Theory)}}\label{i-uxc758uxbbf8uxb860uxc801-uxc774uxb860uxc73cuxb85cuxc11cuxc758-uxc2e4uxc99duxc8fcuxc758-positivism-as-a-semantic-theory}

내 책은 드워킨에 의해 현대 법실증주의의 대표작으로 취급된다. 이 현대
법실증주의는 주로 벤담과 오스틴의 명령 이론들(imperative theories of
law)과, 모든 법이 법적으로 무제한적인 주권적 입법 \textbf{(p.~245)} 인격
또는 기관으로부터 나온다는 그들의 구상을 거부한다는 점에서 그들 같은
초기 판본들과 구별된다. 드워킨은 나의 법실증주의 판본에서 서로 다르지만
관련된 많은 오류들을 발견한다. 이 오류들 가운데 가장 근본적인 것은, 법적
권리들과 법적 책무들(legal duties)을 서술하는 법명제들(propositions of
law) 같은 것들의 참은 개인적 믿음들과 사회적 태도들에 관한 사실들을
포함한 명백한 역사적 사실에 관한 질문들에만 의존한다는
견해이다.\footnote{\emph{LE} 6면 이하.} 법명제들의 참이 의존하는
사실들은 드워킨이 ‘법의 근거들(grounds of law)’이라고 부르는 것을
구성하며,\footnote{\emph{LE} 4면.} 그에 따르면 실증주의자는 이것들이
판사들과 법률가들이 공유하는 언어 규칙들에 의해 고정되어 있다고 잘못
받아들인다. 이 언어 규칙들은 `law'라는 단어가 특정 체계의 특정 쟁점에
관한 ‘당해 법(the law)’이 무엇인지에 대한 진술들에 나타날 때와
‘법(law)’, 곧 일반적 법이 무엇인지에 관한 진술들에 나타날 때 모두 그
단어의 사용, 따라서 의미를 지배한다고 한다.\footnote{\emph{LE} 31면
  이하.} 이러한 실증주의적 법관에 따르면, 법문제들(questions of law)에
관해 있을 수 있는 유일한 의견불일치들은 그러한 역사적 사실들의 존재 또는
부존재에 관한 것들이며, 무엇이 법의 ‘근거들’을 구성하는지에 관한 이론적
의견불일치나 논쟁은 있을 수 없게 된다.

드워킨은 법실증주의 비판의 많은 해명적 지면을, 무엇이 법의 근거들을
구성하는지에 관한 이론적 의견불일치가 실증주의자의 견해와 반대로 영미법
실천들의 두드러진 특징임을 보이는 데 할애한다. 법률가들과 판사들이
공유하는 언어 규칙들에 의해 그것들이 논쟁의 여지 없이 고정되어 있다는
견해에 맞서, 드워킨은 그것들이 본질적으로 논쟁적이라고 주장한다. 그
가운데에는 역사적 사실들뿐 아니라 매우 자주 논쟁적인 도덕 판단들과 가치
판단들도 있기 때문이다.

드워킨은 나 같은 실증주의자들이 어떻게 이 근본적으로 잘못된 견해를
채택하게 되었는지에 관한 매우 다른 두 설명을 제시한다. 이 설명들 가운데
첫째에 따르면, 실증주의자들은 법의 근거들이 규칙들에 의해 논쟁의 여지
없이 고정되어 있지 않고 이론적 의견불일치를 허용하는 논쟁적 사안이라면,
`law'라는 단어가 서로 다른 사람들에게 서로 다른 것을 \emphksans{뜻하게(mean)}
될 것이고, 그 단어를 사용하면서 그들은 같은 것에 관해 소통하는 것이
아니라 서로 빗나가 말하게 될 것이라고 믿는다. 이렇게 실증주의자에게
귀속된 믿음은 드워킨의 견해로는 전적으로 잘못된 것이며, 그는
\textbf{(p.~246)} 실증주의자가 그것에 기초한다고 상정되는 논쟁적 법의
근거들에 대한 논변을 ‘의미론적 독침(semantic sting)’\footnote{\emph{LE}
  45면.}이라고 부른다. 그것이 'law'라는 단어의 의미에 관한 이론에
기초하기 때문이다. 그래서 그는 \emph{Law's Empire}에서 이 ‘의미론적
독침’을 뽑아내고자 했다.

\emph{Law's Empire}의 첫 장에서 나는 오스틴과 함께 의미론적 이론가로
분류되고, 그래서 'law'라는 단어의 의미에서 명백한-사실 실증주의적
법이론을 도출하며 의미론적 독침을 겪는 것으로 분류된다. 그러나 실제로 내
책이나 내가 쓴 다른 어떤 것에도 내 이론에 관한 그러한 설명을 뒷받침하는
것은 없다. 따라서 발전된 국내법체계들이 법원들이 적용해야 하는 법들의
식별 기준들을 특정하는 승인 규칙을 포함한다는 나의 교설은 틀릴 수
있지만, 나는 어디에서도 이 교설을 모든 법체계들 안에는 그러한 승인
규칙이 있어야 한다는 것이 'law'라는 단어의 의미의 일부라는 잘못된 관념에
기초시키지 않았고, 법의 근거들의 식별 기준들이 논쟁의 여지 없이 고정되어
있지 않다면 'law'가 서로 다른 사람들에게 서로 다른 것을 \emphksans{뜻하게} 될
것이라는 훨씬 더 잘못된 관념에 기초시키지도 않았다.

기실 내게 귀속된 이 마지막 논변은 한 개념의 \emphksans{의미}와 그 개념의
\emphksans{적용} 기준들을 혼동한다. 나는 이것을 받아들이기는커녕
정의(justice)의 개념을 설명하면서 이 책 160쪽에서, 일정한 의미를 가진
개념의 적용 기준들이 달라질 수도 있고 논쟁적일 수도 있다는 사실에
명시적으로 주의를 환기했다. 이것을 명료하게 하기 위해 나는 사실상
드워킨의 후기 작업에서 그토록 두드러지게 나타나는 개념(concept)과 한
개념의 서로 다른 구상들(conceptions) 사이의 바로 그 구별을
그었다.\footnote{이 구분에 대해서는 John Rawls, \emph{A Theory of
  Justice} (1971), 5--6, 10면 참조. {[}Rawls는 정의 개념(concept of
  justice)과 정의의 구상들(conceptions of justice)을 구분하면서 “나는
  H. L. A. Hart의 \emph{The Concept of Law}\ldots{} 155--159면을
  따른다”고 명시함. \emph{A Theory of Justice} 5면 각주 1 참조.{]}}

마지막으로 드워킨은 또한, 실증주의자가 자신의 법이론이 의미론적 이론이
아니라 복합적 사회현상으로서 일반적 법의 독특한 특징들에 관한 기술적
설명이라고 주장하는 것은, 의미론적 이론과의 대비를 제시하지만 그 대비는
공허하고 오도적이라고 고집한다. 그의 논변\footnote{\emph{LE} 418--419면,
  주석 29.}은, 사회현상으로서 법의 독특한 특징들 가운데 하나가
법률가들이 법명제들의 참을 논쟁하고 \textbf{(p.~247)} 이것을 그러한
명제들의 의미를 참조하여 ‘설명한다’는 것이므로, 그러한 기술적 법이론은
결국 의미론적이어야 한다는 것이다.\footnote{\emph{LE} 31--33면 참조.} 이
논변은 내게는 'law'의 의미와 법명제들의 의미를 혼동하는 것으로 보인다.
법에 관한 의미론적 이론은, 'law'라는 단어의 바로 그 의미가 법을 일정한
특정 기준들에 의존하게 만든다는 이론이라고 드워킨에 의해 말해진다.
그러나 법명제들은 전형적으로 'law'가 무엇인지가 아니라 \emphksans{당해 법(the
law)}이 무엇인지, 곧 어떤 체계의 법이 사람들에게 무엇을 하도록
허용하거나, 요구하거나, 권한을 부여하는지에 관한 진술들이다. 그러므로
그러한 법명제들의 의미가 정의들이나 진리조건들에 의해 결정된다 하더라도,
이것은 'law'라는 단어의 바로 그 의미가 법을 일정한 특정 기준들에
의존하게 만든다는 결론으로 이어지지 않는다. 그러한 경우가 되려면 한
체계의 승인 규칙이 제공하는 기준들과 그러한 규칙의 필요성이 'law'라는
단어의 의미에서 도출되어야 할 것이다. 그러나 내 작업에는 그러한 교설의
흔적이 전혀 없다.\footnote{본서 209면 참조. 이곳에서 나는 그러한 교설을
  명시적으로 거부하고 있다.}

드워킨이 나의 법실증주의 형식을 잘못 제시하는 또 하나의 측면이 있다.
그는 나의 승인 규칙 교설을, 승인 규칙이 법의 식별을 위해 제공하는
기준들이 오직 역사적 사실들로만 이루어져야 한다고 요구하는 것으로,
그래서 ‘명백한-사실 실증주의(plain-fact positivism)’의 한 사례로
취급한다.\footnote{{[}이 표현은 하트 자신의 것이며, \emph{LE}에는
  등장하지 않음.{]}} 그러나 승인 규칙이 제공하는 기준들에 관해 내가 든
주된 예들이, 드워킨이 ‘계보(pedigree)’라고 부른 사항들, 곧 법들이
법제도들에 의해 채택되거나 창설되는 방식에만 관련되고 그 내용에는
관련되지 않는 사항들인 것은 사실이지만,\footnote{\emph{TRS} 17면.} 나는
이 책(72쪽)과 앞선 논문 「실증주의와 법과 도덕의 분리(Positivism and the
Separation of Law and Morals)」 양쪽에서, 미국에서처럼 어떤
법체계들에서는 법적 유효성의 궁극적 기준들이 계보뿐 아니라 정의
원칙들이나 실체적 도덕 가치들(substantive moral values)을 명시적으로
편입할 수도 있고, 이것들이 법적 헌법적 제약들의 내용을 형성할 수도
있다고 명시적으로 진술했다.\footnote{\emph{Harvard Law Review} 제71권
  (1958) 598면. \emph{Essays on Jurisprudence and Philosophy}에도
  재수록. 특히 54--55면 참조.} \emph{Law's Empire}에서 드워킨이 내게
‘명백한-사실’ 실증주의를 귀속시키는 것은 내 이론의 이 측면을 무시하는
것이다. 그러므로 \textbf{(p.~248)} 그가 내게 귀속시키는 명백한-사실
실증주의의 ‘의미론적’ 판본은 분명 내 것이 아니며, 내 이론은 어떤 형식의
명백한-사실 실증주의도 아니다.

\subsubsection{\texorpdfstring{(ii) \emphksans{해석적 이론으로서의 실증주의}
(Positivism as an Interpretive
Theory)}{(ii) 해석적 이론으로서의 실증주의 (Positivism as an Interpretive Theory)}}\label{ii-uxd574uxc11duxc801-uxc774uxb860uxc73cuxb85cuxc11cuxc758-uxc2e4uxc99duxc8fcuxc758-positivism-as-an-interpretive-theory}

드워킨의 명백한-사실 실증주의에 대한 둘째 설명은 그것을 의미론적
이론이나 언어적 고려들에 기초한 것으로 다루지 않고, 그가
‘관행주의(conventionalism)’라고 부르는 드워킨식 해석이론의 한 형식으로
재구성하려고 시도한다. 이 이론, 곧 드워킨이 결국 결함 있는 것으로
거부하는 이론에 따르면, 실증주의자는 법을 최선의 빛 아래 보이는 데
전념하는 해석 이론가라는 외양 아래에서, 법의 기준들을 명백한 사실들로
이루어진 것으로 제시한다. 다만 이 사실들은 의미론적 판본에서처럼 법의
어휘에 의해 고정되는 것이 아니라, 판사들과 법률가들이 공유하는 확신에
의해 논쟁의 여지 없이 고정된다. 이것은 법을 수범자들에게 큰 가치가 있는
어떤 것을 확보해 주는 것으로 보임으로써 법에 유리한 빛을 던진다. 곧 법적
강제의 계기들이 모두에게 이용 가능한 명백한 사실들에 의존하도록 만들어,
강제가 사용되기 전에 모두가 공정한 경고를 받을 수 있게 한다는 것이다.
드워킨은 이것을 ‘보호된 기대들의 이상(the ideal of protected
expectations)’이라고 부르지만,\footnote{\emph{LE} 117면.} 그에게 그
이점들은 결국 그 여러 결함들을 능가하지 않는다.

그러나 관행주의로서의 실증주의에 관한 이 해석주의적 설명은 내 법이론의
그럴듯한 판본이나 재구성으로 제시될 수 없다. 두 가지 이유에서 그렇다.
첫째, 내가 이미 말했듯이 내 이론은 명백한-사실 실증주의 이론이 아니다.
법의 기준들 가운데 ‘명백한’ 사실들뿐 아니라 가치들도 인정하기 때문이다.
둘째, 그리고 더 중요하게, 드워킨의 해석적 법이론은 그 모든 형식에서 법과
법적 실천의 요점 또는 목적이 강제를 정당화하는 것이라는 전제에 기초해
있지만,\footnote{{[}\emph{LE} 93면.{]}} 법이 이것을 자신의 요점 또는
목적으로 가진다는 것은 결코 내 견해가 아니며 그런 적도 없다. 다른
실증주의 형식들처럼 내 이론은 법과 법적 실천들 그 자체의 요점 또는
목적을 식별한다고 주장하지 않는다. 그러므로 법의 목적이 강제 사용을
정당화하는 것이라는 드워킨의 견해, 내가 결코 공유하지 않는 그 견해를
뒷받침하는 것은 내 이론 안에 없다. \textbf{(p.~249)} 사실 나는 법 그
자체가 봉사하는 더 특정한 목적을, 인간 행위에 대한 지침들과 그러한
행위에 대한 비판 기준들을 제공하는 것 너머에서 찾는 것은 매우 헛되다고
생각한다. 물론 이것은 같은 일반적 목적들을 가진 다른 규칙들이나
원칙들로부터 법들을 구별하는 데에는 충분하지 않을 것이다. 법의 독특한
특징들은 그 기준들의 식별, 변경, 집행을 위해 이차 규칙들을 통해 마련하는
장치와, 다른 기준들에 대한 우선성(priority)을 내세우는 일반적 주장에
있다. 그러나 내 이론이 법적 강제의 계기들에 대한 사전 고지가 일반적으로
이용 가능하리라는 점을 보증함으로써 기대들을 보호하는 관행주의 형식의
명백한-사실 실증주의에 전적으로 결속되어 있다 하더라도, 이것은 오직 내가
이것을 법이 가진 특정한 도덕적 이점(moral merit)으로 본다는 것만을 보여
줄 뿐, 법 그 자체의 전체 목적이 이것을 제공하는 것임을 보여 주지는 않을
것이다. 법적 강제의 계기들은 주로 그 수범자들의 행위를 인도하는 법의
일차적 기능이 무너진 경우들이므로, 법적 강제는 물론 중요한 사안이지만
부차적 기능이다. 그것의 정당화는 법 그 자체의 요점 또는 목적으로
온당하게 취급될 수 없다.

법적 강제는 그것이 “관행적 이해들(conventional understandings)에 부합할
때”\footnote{\emph{LE} 429면 주석 3.}에만 정당화된다는 주장을 하는
관행주의적 해석이론으로 내 법이론을 재구성하려는 드워킨의 이유들은, 이
책 제V장 제3절의 ‘법의 요소들(Elements of Law)’에 관한 내 설명에
기초한다. 거기서 나는 승인, 변경, 재판의 이차 규칙들을, 오직 의무의 일차
규칙들만으로 이루어진 상상된 단순한 체제의 결함들에 대한
치유책들(remedies)로 제시한다. 이 결함들은 규칙들의 동일성에 관한
\emphksans{불확실성}, 규칙들의 \emphksans{정태적} 성질, 그리고 규칙들이 오직
분산된 사회적 압력에 의해서만 집행될 때의 시간낭비적
\emphksans{비효율성}이다. 그러나 이러한 이차 규칙들을 그러한 결함들에 대한
치유책들로 제시하면서도, 나는 어디에서도 법적 강제가 이러한 규칙들에
부합할 때에만 \emphksans{정당화된다}고 주장하지 않았고, 더더욱 그러한
정당화의 제공이 일반적 법의 요점 또는 목적이라고 주장하지 않았다. 기실
이차 규칙들에 관한 논의에서 내가 강제를 언급하는 유일한 대목은 규칙들의
집행을 \textbf{(p.~250)} 법원들이 운용하는 조직화된 제재들 대신 분산된
사회적 압력에 맡기는 일의 시간낭비적 \emphksans{비효율성}에 관한 것이다.
그러나 분명 비효율성에 대한 치유책은 정당화가 아니다.

물론 의무의 일차 규칙들의 체제에 이차적 승인 규칙이 추가되면, 개인들이
강제의 계기들을 사전에 식별할 수 있게 하는 경우가 많으므로, 강제 사용에
대한 하나의 도덕적 반론을 배제한다는 의미에서 그 사용을 정당화하는 데
도움이 될 것이다. 그러나 승인 규칙이 가져올 법의 요구사항들에 관한
확실성과 사전 지식은 강제가 문제되는 경우에만 중요한 것이 아니다. 그것은
예컨대 유언장을 작성하거나 계약을 체결할 행위능력으로 나타나는 법적
권한들의 지성적 행사와, 일반적으로 사적·공적 삶의 지성적 계획에도 똑같이
핵심적이다. 그러므로 승인 규칙이 기여하는 강제의 정당화는 그 승인 규칙의
일반적 요점 또는 목적으로 제시될 수 없고, 더더구나 법 전체의 일반적 요점
또는 목적으로 제시될 수 없다. 내 이론에는 그것이 그럴 수 있다고 시사하는
것이 없다.

\subsubsection{\texorpdfstring{(iii) \emphksans{연성 실증주의} (Soft
Positivism)}{(iii) 연성 실증주의 (Soft Positivism)}}\label{iii-uxc5f0uxc131-uxc2e4uxc99duxc8fcuxc758-soft-positivism}

드워킨은 내게 ‘명백한-사실 실증주의’ 교설을 귀속시키면서, 내 이론을 승인
규칙의 존재와 권위가 법원들에 의한 그 수용이라는 사실에 의존해야 한다고
요구하는 것으로는, 실제로 그렇듯이, 취급했을 뿐 아니라, 그 규칙이
제공하는 법적 유효성 기준들이 그가 ‘계보(pedigree)’ 사항들이라고 부르는,
법의 창설 또는 채택의 방식과 형식에 관한 특정한 종류의 명백한 사실로만
이루어져야 한다고 요구하는 것으로도, 실제로는 그렇지 않은데, 취급하였다.
이것은 이중으로 잘못되었다. 첫째, 그것은 승인 규칙이 법적 유효성
기준들로서 도덕 원칙들 또는 실체적 가치들(substantive values)에 대한
부합(conformity)을 편입할 수 있다는 나의 명시적 인정을 무시한다. 따라서
내 교설은 이른바 ‘연성 실증주의(soft positivism)’이지, 드워킨의 그
판본에서처럼 ‘명백한-사실’ 실증주의가 아니다. 둘째, 승인 규칙이 제공하는
명백한-사실 기준들이 오직 계보 사항들일 수밖에 없다고 시사하는 것은 내
책에 전혀 없다. 그것들은 대신 종교의 설립이나 투표권 제한에 관한 미국
헌법 제16 수정조항 또는 제19 수정조항처럼 입법 내용에 대한 실체적
제약들(substantive constraints)일 수 있다.

그러나 이 답변은 드워킨의 가장 기본적인 비판들에 답하지 못한다.
\textbf{(p.~251)} 연성 실증주의의 어떤 형식을 채택한 다른 이론가들에게
답하면서,\footnote{E. P. Soper와 J. L. Coleman에 대한 드워킨의 응답을
  \emph{RDCJ} 247면 이하, 252면 이하에서 볼 수 있음.} 그는 그것이
타당하다면 내 이론에도 적용되어 여기서 답을 요구할 중요한 비판들을
제기했기 때문이다.

드워킨의 가장 근본적인 비판은, 법의 식별이 도덕 판단 또는 다른 가치
판단에 대한 부합이라는 논쟁적 사안들에 의존할 수 있도록 허용하는 연성
실증주의와, 논쟁적 도덕 논변들에 의존하지 않고 명백한 사실의 문제들로서
확실하게 식별될 수 있는 신뢰할 만한 공적 행위 기준들을 제공하는 데 법이
본질적으로 관련되어 있다는 일반 실증주의적 ‘그림(picture)’ 사이에 깊은
비일관성이 있다는 것이다.\footnote{\emph{RDCJ} 248면.} 연성 실증주의와
내 이론의 나머지 부분 사이의 그러한 비일관성을 입증하기 위해 드워킨은,
다른 결함들 가운데서도 상상된 전법적(pre-legal) 관습-유형 의무 일차 규칙
체제의 불확실성을 치유하는 것으로 승인 규칙을 설명한 내 설명을 인용할
것이다. 연성 실증주의에 관한 이 비판은 내게는 일관된 실증주의자가 법적
기준들의 한 집합에 귀속해야 한다고 보는 확실성의 정도와, 법적 유효성
기준들이 특정 도덕 원칙들이나 가치들에 대한 부합을 포함할 경우 생길
불확실성 양쪽 모두를 과장하는 것으로 보인다. 승인 규칙의 중요한 기능
하나가 법이 확인될 수 있는 확실성을 증진하는 것임은 물론 참이다. 승인
규칙이 법에 대해 도입한 검사 기준들이 어떤 사례들에서 논쟁적 쟁점들을
제기할 뿐 아니라 모든 사례 또는 대부분의 사례에서 그렇게 한다면, 그것은
이 기능을 수행하지 못할 것이다. 그러나 다른 가치들의 측면에서 어떤
대가를 치르더라도 모든 불확실성을 배제하는 것은 내가 승인 규칙을 위해
생각해 본 목표가 아니다. 이는, 승인 규칙 자체뿐 아니라 그것을 참조하여
식별되는 특정 법규칙들도 다툴 수 있는 불확실성의 ‘반음영(penumbra)’을
가질 수 있다는 내 명시적 진술에 의해서도, 내가 바랐듯이,
분명해진다.\footnote{본서 123면, 147--154면 참조.} 또한 법들이 그 의미에
관해 생길 수 있는 모든 가능한 질문들을 미리 확정할 수 있도록 틀 지어질
수 있다 하더라도, 그러한 법들을 채택하는 것은 법이 소중히 여겨야 할 다른
목적들과 자주 충돌할 것이라는 나의 일반 논변도 있다.\footnote{{[}본서
  128면 참조.{]}} 예견되지 않은 사례의 구성이 알려지고 그 결정에서 걸린
쟁점들이 식별되어 합리적으로 확정될 수 있도록, 많은 법규칙들의 경우
일정한 불확실성의 여지는 \textbf{(p.~252)} 용인되어야 하며 기실
환영되어야 한다. 승인 규칙의 확실성-제공 기능을 최우선적이고 압도적인
것으로 취급할 때에만, 논쟁적일 수 있는 도덕 원칙들이나 가치들에 대한
부합을 법의 기준들 가운데 포함하는 연성 실증주의의 형식이 비일관적인
것으로 간주될 수 있을 것이다. 여기서 바탕에 있는 질문은, 법체계가
일반적으로 신뢰할 수 있고 사전에 식별 가능한 확정적 행위 지침들을
제공함으로써 탈중심화된 관습-유형 규칙 체제로부터 어떤 유의미한 진전을
이루려면, 어느 정도 또는 범위의 불확실성을 용인할 수 있는가에 관한
것이다.

나의 연성 실증주의 판본의 일관성에 대한 드워킨의 둘째 비판은 법의
확정성(determinacy)과 완전성(completeness)에 관한 서로 다르고 더 복잡한
쟁점들을 제기한다. 이 책에서 전개된 내 견해는, 승인 규칙이 제공하는
기준들에 의해 일반적 항들로 식별되는 법규칙들과 법원칙들이 내가 자주
‘개방적 직조(open texture)’라고 부르는 것을 종종 가진다는 것이다. 그래서
주어진 규칙이 특정 사례에 적용되는지가 문제될 때, 법은 어느 쪽의 답도
확정하지 못하고 따라서 부분적으로 불확정적인 것으로 드러난다. 그러한
사례들은 단지 합리적이고 식견 있는 법률가들이 어느 답이 법적으로
올바른지에 관하여 의견불일치할 수 있다는 의미에서 논쟁적인 ‘어려운
사례들’일 뿐 아니라, 그러한 경우 법은 근본적으로 \emphksans{불완전하다}.
그것은 그러한 사례들에서 쟁점이 된 질문들에 \emphksans{아무} 답도 제공하지
않는다. 그것들은 법적으로 규율되지 않으며, 그러한 사례들에서 결정에
이르기 위해 법원들은 내가 ‘재량(discretion)’이라고 부르는 제한된 법창출
기능을 행사해야 한다. 드워킨은 법이 이런 방식으로 불완전할 수 있고
그러한 법창출 재량의 행사로 채워질 틈들을 남길 수 있다는 관념을
거부한다. 그는 이 견해가, 법적 권리 또는 법적 책무의 존재를 단언하는
법명제가 논쟁적일 수 있고 따라서 합리적이고 식견 있는 사람들이 의견을
달리할 수 있으며, 그들이 의견을 달리할 때 그것이 참인지 거짓인지를
확정적으로 논증할 방법이 자주 없다는 사실에서 나온 잘못된 추론이라고
생각한다. 그러한 추론이 잘못인 것은, 법명제가 이렇게 논쟁적일 때에도
그럼에도 그것을 참 또는 거짓으로 만드는 ‘사안의 사실들(facts of the
matter)’이 있을 수 있고, \textbf{(p.~253)} 그 참이나 거짓을 논증할 수
없더라도 그것이 참이라는 논변들이 그것이 거짓이라는 논변들보다 더 나은
것으로 평가될 수 있으며, 그 반대도 가능하기 때문이다. 논쟁적인 법과
불완전하거나 불확정적인 법 사이의 이 구별은 드워킨의 해석이론에 상당히
중요하다. 그 이론에 따르면 법명제는 다른 전제명제들과 결합하여, 법체계의
제도적 역사에 가장 잘 부합하고 그것에 대한 최선의 도덕적 정당화도
제공하는 원칙들에서 따라 나올 때에만 참이기 때문이다. 그러므로
드워킨에게 어떤 법명제의 참은 궁극적으로 무엇이 가장 잘 정당화하는지에
관한 도덕 판단의 참에 의존하며, 그에게 도덕 판단들은 본질적으로
논쟁적이므로 모든 법명제들도 그러하다.

드워킨에게, 적용이 논쟁적인 도덕 판단을 포함하는 법적 유효성 기준이라는
관념은 이론적 어려움을 전혀 제시하지 않는다. 그의 견해에서는 그것이
여전히 이미 존재하는 법(pre-existing law)에 대한 진정한 검사
기준(test)일 수 있다. 그 논쟁적 성격은 그것을 참으로 만드는 사실들, 많은
경우 도덕적 사실들이 존재한다는 점과 완전히 양립 가능하기 때문이다.

그러나 법적 유효성 기준이 부분적으로 도덕적 검사 기준일 수 있음을
허용하는 연성 실증주의는, 드워킨의 주장에 따르면, 위 pp.~251-252에서
이미 논의한 것에 더해 둘째 비일관성에 연루된다. 그것은 법이 확실하게
식별 가능하다는 실증주의적 ‘그림’과 비일관적일 뿐 아니라, 법명제들의
‘객관적 지위(objective standing)’\footnote{\emph{RDCJ} 250면.}를 도덕
판단들의 지위에 관한 어떤 논쟁적 철학 이론에 대한 결속으로부터도
독립시키려는, 그가 실증주의자들에게 귀속시키는 바람과도 비일관적이다.
도덕적 검사 기준이 이미 존재하는 법에 대한 검사 기준일 수 있으려면, 도덕
판단들이 참이 되는 객관적 도덕 사실들이 있어야 한다. 그러나 그러한
객관적 도덕 사실들이 있다는 것은 논쟁적 철학 이론이다. 그러한 사실들이
없다면, 도덕적 검사 기준을 적용하라는 말을 들은 판사는 이것을 오직,
법체계가 이에 부과하는 어떤 제약들에 종속하여, 도덕과 그 요구사항들에
관한 자신의 최선의 이해에 따라 법창출 재량을 행사하라는 요청으로만
취급할 수 있다.

나는 여전히 법이론이 도덕 판단들의 일반적 지위에 관한 \textbf{(p.~254)}
논쟁적 철학 이론들에 결속되는 것을 피해야 하며, 이 책에서 내가
그러하듯(p.~168), 그것들이 드워킨이 ‘객관적 지위’라고 부르는 것을
가지는지에 관한 일반 질문을 열어 두어야 한다고 생각한다. 이 철학적
질문에 대한 답이 무엇이든, 판사의 책무는 동일할 것이기 때문이다. 곧 그가
결정해야 할 수 있는 어떤 도덕적 쟁점들에 대해서도 자신이 할 수 있는
최선의 도덕 판단을 내리는 것이다. 사건들을 이렇게 결정하면서 판사가
도덕에 따라, 법이 부과하는 어떤 제약들에 종속하여, 법을 \emphksans{만들고}
있는 것인지, 아니면 법에 대한 도덕적 검사 기준에 의해 드러나는 이미
\emphksans{존재하는} 법이 무엇인지에 관한 자신의 도덕 판단에 의해 인도되는
것인지는 어떤 실천적 목적에서도 중요하지 않을 것이다. 물론, 내가 그래야
한다고 주장하듯이 법이론이 도덕 판단들의 객관적 지위 문제를 열어 둔다면,
연성 실증주의는 단순히 도덕 원칙들이나 가치들이 법적 유효성의 기준들
가운데 있을 수 있다는 이론으로 특징지어질 수 없다. 도덕 원칙들과
가치들이 객관적 지위를 가지는지가 열린 질문이라면, 그것들에 대한 부합을
현존 법(existing law)에 대한 검사 기준들 가운데 포함하려는 ‘연성
실증주의적’ 규정들이 그러한 효과를 가질 수 있는지, 아니면 오직
법원들에게 도덕에 따라 법을 \emphksans{만들라}는 지시들을 구성할 수 있을
뿐인지도 열린 질문이어야 하기 때문이다.

몇몇 이론가들, 특히 라즈가 다음과 같이 본다는 점은 관찰되어야 한다. 도덕
판단들의 지위가 무엇이든, 법이 법원들에게 법이 무엇인지를 확정하는 데
도덕 기준들을 적용하라고 요구할 때마다, 법은 그로써 법원들에게 재량을
부여하고, 새 법을 만들면서 자기들의 최선의 도덕 판단에 따라 그 재량을
사용하라고 지시한다. 법은 그로써 도덕을 이미 존재하는 법으로 전환하지
않는다.\footnote{J. Raz, “Dworkin: A New Link in the Chain”,
  \emph{California Law Review} 제74권 (1986) 1103면, 1110면 및
  1115--1116면.}

\subsection{3. 규칙의 본성 (THE NATURE OF
RULES)}\label{uxaddcuxce59uxc758-uxbcf8uxc131-the-nature-of-rules}

\subsubsection{\texorpdfstring{(i) \emphksans{규칙의 실천 이론} (The Practice
Theory of
Rules)}{(i) 규칙의 실천 이론 (The Practice Theory of Rules)}}\label{i-uxaddcuxce59uxc758-uxc2e4uxcc9c-uxc774uxb860-the-practice-theory-of-rules}

이 책의 여러 지점에서 나는 법의 내부 진술(internal statements)과 외부
진술(external statements) 사이의 구분, 그리고 법의 내부 측면(internal
aspects)과 외부 측면(external aspects) 사이의 구분에 주의를 기울였다.

이러한 구분들과 그 중요성을 설명하기 위해 나는 (56--7쪽)
\textbf{(p.~255)} 제정된 규칙들과 관습-유형 규칙들 양쪽을 포함하는
고도로 복잡한 법체계의 사례에서 출발하지 않고, 그보다 단순한 사례, 곧
크든 작든 어떤 사회 집단의 관습-유형 규칙들의 사례에서 출발했다.
내부/외부 구분은 이 사례에도 똑같이 적용되며, 나는 이러한 규칙들을
‘사회적 규칙들(social rules)’이라고 부른다. 내가 그것들에 관해 제시한
설명은 ‘규칙의 실천 이론(the practice theory)’으로 알려지게 되었는데,
그것은 한 집단의 사회적 규칙들을 사회적 실천의 한 형식에 의해 구성되는
것으로 취급하기 때문이다. 이 사회적 실천은 집단 구성원 대부분이
반복적으로 일양적인 따름을 보이는 행위 양식들과, 그러한 행위 양식들에
대한 독특한 규범적 태도, 곧 내가 ‘수용(acceptance)’이라고 부른 것을 함께
포함한다. 이 수용은 그러한 행위 양식들을 자신의 미래 행위의 지침들로,
그리고 부합(conformity)에 대한 요구들과 여러 형태의 압력을 정당화할 수
있는 비판 기준들로 삼으려는 개인들의 지속적 성향으로 이루어진다. 사회적
규칙들에 대한 외부 관점은 그 실천의 관찰자의 관점이고, 내부 관점은 그
실천의 참여자, 곧 규칙들을 행위의 지침들과 비판 기준들로 수용하는
참여자의 관점이다.

나의 사회적 규칙들의 실천 이론은 드워킨에 의해 광범위하게 비판되어 왔다.
앞서 말했듯이 그는 이와 유사하지만 실제로는 여러 점에서 매우 다른 구별,
곧 한 공동체의 사회적 규칙들에 대한 사회학자의 외부 기술과, 자신의
행위와 타인의 행위에 대한 평가와 비판을 위해 규칙들에 호소하는 참여자의
내부 관점 사이의 구별을 제시한다.\footnote{{[}\emph{LE} 13--14면
  참조.{]}} 사회적 규칙들에 대한 나의 원래 설명에 관해 드워킨이 제기한
비판 가운데 일부는 분명 타당하며 법의 이해에도 중요하다. 이어서 나는
이제 필요하다고 보는 나의 원래 설명의 상당한 수정들을 지목한다.

\begin{enumerate}
\def\labelenumi{(\roman{enumi})}
\item
  드워킨이 주장했듯이, 나의 설명은 한 집단의 관행적 규칙들에서 드러나는
  \emphksans{관행(convention)}의 합의와, 한 집단의 동시적 실천들에서 드러나는
  독립적 \emphksans{확신(conviction)}들의 합의 사이의 중요한 차이를
  무시한다는 점에서 결함이 있다. 어떤 규칙들에 대한 집단의 일반적
  부합(conformity)이 그 개별 구성원들이 그 규칙들을 수용하는 이유들의
  일부라면, 그 규칙들은 관행적 사회 실천들이다. 이와 대조적으로
  \textbf{(p.~256)} 한 집단의 공유 도덕과 같은 단지 동시적인 실천들은
  관행이 아니라, 집단 구성원들이 특정 방식으로 행위할 동일하지만
  독립적인 이유들을 가지고 대체로 그 이유들에 따라 행위한다는 사실에
  의해 구성된다.
\item
  드워킨이 또한 옳게 주장했듯이, 나의 사회적 규칙 설명은 내가 이제
  설명한 의미에서 관행적인 규칙들에만 적용 가능하다. 이것은 나의 실천
  이론의 범위를 상당히 좁히며, 나는 이제 그것을 개인적 도덕이든 사회적
  도덕이든 도덕에 대한 건전한 설명으로 보지 않는다. 그러나 그 이론은
  일반적 사회 관습들, 곧 법적 효력을 가진 것으로 승인될 수도 있고 그렇지
  않을 수도 있는 관습들뿐 아니라, 승인 규칙을 포함한 일정한 중요한
  법규칙들도 포함하는 관행적 사회적 규칙들에 대한 충실한 설명으로는 남아
  있다. 승인 규칙은 사실상 법원들의 법-식별 및 법-적용 작용에서 수용되고
  실천될 때에만 존재하는 사법적 관습 규칙의 한 형식이다. 이와 대조적으로
  제정된 법규칙들은, 승인 규칙이 제공하는 기준들에 의해 유효한
  법규칙들로 식별될 수 있지만, 그것들을 실천할 계기가 생기기 전에도
  제정의 순간부터 법규칙들로 존재할 수 있으므로, 실천 이론은 그것들에는
  적용되지 않는다.
\end{enumerate}

드워킨의 규칙의 실천 이론에 대한 핵심 비판은, 그것이 사회적 규칙을 그
사회적 실천에 의해 구성되는 것으로 잘못 받아들이며, 그래서 그러한 규칙이
존재한다는 진술을 단지 그 규칙의 존재를 위한 실천-조건들이 충족되었다는
외부 사회학적 사실의 진술로 취급한다는 것이다.\footnote{{[}\emph{TRS}
  48--58면 참조.{]}} 드워킨의 논변에 따르면 그러한 설명은 가장 단순한
관행적 규칙조차 가지는 \emphksans{규범적} 성격을 설명할 수 없다. 그러한
규칙들은 \emphksans{책무들}과 \emphksans{행위의 이유들}을 확립하며, 그러한
규칙들이 행위 비판과 행위 요구의 뒷받침에서 통상 인용될 때에는 바로
이것들에 호소하기 때문이다. 규칙들의 이러한 이유-제공 및 책무-확립
특징이 그 독특한 규범적 성격을 구성하며, 그것은 규칙들의 존재가 실천
이론에 따르면 사회적 규칙의 존재를 구성하는 실천들과 태도들이 그러하듯이
단지 사실적 사태로 이루어질 수 없음을 보여 준다. 드워킨에 따르면, 이러한
독특한 특징들을 가진 규범적 규칙은 ‘일정한 규범적 사태’가 있을 때에만
존재할 수 있다.\footnote{\emph{TRS} 51면.} \textbf{(p.~257)} 나는 인용된
이 말들이 답답할 만큼 모호하다고 본다. ‘교회 출석자들의 규칙’(남성은
교회에서 머리를 드러내야 한다)의 예에 관한 논의에서,\footnote{{[}\emph{TRS}
  50--58면; 본서 124--125면 참조.{]}} 드워킨은 규범적 사태로 규칙이
요구하는 것을 할 좋은 도덕적 근거들 또는 정당화의 존재를 뜻하는 듯하다.
그래서 그는 교회 출석자들이 교회에서 모자를 벗는 단순한 반복적으로
일양적인 실천은 그 규칙을 구성할 수 없지만, 불쾌감을 주는 방식들을
창출하고 교회에서 모자를 벗도록 요구하는 규칙을 위한 좋은 근거들인
기대들을 일으킴으로써 그 규칙을 정당화하는 데 도움을 줄 수 있다고
논변한다. 규범적 규칙의 단언을 보증하기 위해 요구되는 규범적 사태로
드워킨이 뜻하는 것이 이것이라면, 사회적 규칙의 존재 조건에 관한 그의
설명은 내게 지나치게 강해 보인다. 그것은 규칙들에 책무를 확립하거나
행위의 이유들을 제공하는 것으로 호소하는 참여자들이 규칙들에 부합할 좋은
도덕적 근거들 또는 정당화가 있다고 믿어야 할 뿐 아니라, 실제로 그러한
좋은 근거들이 있어야 한다고 요구하는 것처럼 보이기 때문이다. 분명 어떤
사회는 그 구성원들이 수용하지만 도덕적으로 사악한 규칙들, 예컨대 특정
피부색을 가진 사람들이 공원이나 해수욕장 같은 공공시설을 이용하는 것을
금지하는 규칙들을 가질 수 있다. 기실 사회적 규칙의 존재를 위한 일반
조건으로, 참여자들이 그 규칙에 부합할 좋은 도덕적 근거들이 있다고
\emphksans{믿어야} 한다는 더 약한 조건조차도 지나치게 강하다. 어떤 규칙들은
단순히 전통에 대한 존중, 타인들과 동일시하려는 바람, 또는 사회가
개인들에게 무엇이 유리한지를 가장 잘 안다는 믿음에서 수용될 수 있기
때문이다. 이러한 태도들은 그 규칙들이 도덕적으로 문제 있다는 어느 정도
생생한 자각과 공존할 수 있다. 물론 관행적 규칙은 도덕적으로 건전하고
정당화될 수도 있으며 그렇게 믿어질 수도 있다. 그러나 관행적 규칙들을
자신의 행위 지침이나 비판 기준들로 수용한 사람들이 왜 그렇게 했는지가
문제될 때, 나는 제시될 수많은 답변들(이 책 203쪽, 232쪽 참조) 가운데
규칙들의 도덕적 정당화에 대한 믿음을 유일하게 가능하거나 충분한 답변으로
골라낼 이유가 없다고 본다.

마지막으로, 드워킨은 규칙의 실천 이론이 \textbf{(p.~258)} 관행적
규칙들로 제한되더라도 폐기되어야 한다고 논변한다. 그것은 관행적 규칙의
범위가 논쟁적일 수 있고 따라서 의견불일치의 대상이 될 수 있다는 관념을
수용할 수 없기 때문이다.\footnote{{[}\emph{TRS} 58면.{]}} 그는
반복적으로 일양적인 실천과 수용에 의해 구성되는 논쟁의 여지 없는
규칙들이 일부 있다는 점은 부정하지 않지만, 그렇게 구성되는 규칙들은 어떤
게임들의 규칙들처럼 상대적으로 중요하지 않은 사례들만 포함한다고
주장한다. 그러나 이 책에서는 법체계의 기본 승인 규칙처럼 중요하고도 거의
논쟁의 여지가 없는 규칙이 법원들이 그것을 자신들의 법-적용 및 법-집행
작용의 지침으로 수용하는 일양적 실천에 의해 구성되는 규칙으로 취급된다.
이에 맞서 드워킨은 어려운 사례들에서 어떤 주제에 관한 당해 법이
무엇인지에 관해 판사들 사이에 이론적 의견불일치가 자주 있으며, 이것들이
논쟁 없음과 일반적 수용의 외양이 환상임을 보여 준다고 주장한다. 물론
그러한 의견불일치들의 빈도와 중요성은 부정될 수 없다. 그러나 그것들의
존재에 호소하여 승인 규칙에 실천 이론이 적용될 수 없다고 논변하는 것은
그 규칙의 기능에 대한 오해에 기초한다. 그것은 승인 규칙이 특정 사례들의
법적 결과를 완전히 확정하도록 의도되어, 어떤 사례에서 생기는 어떤 법적
쟁점도 그 규칙이 제공하는 기준들이나 검사 기준들에 대한 단순한
호소만으로 해결될 수 있다고 가정한다. 그러나 이것은 오해이다. 그 규칙의
기능은 현대 법체계들에서 올바른 법적 결정들이 충족해야 하는 일반
조건들만을 확정하는 것이다. 그 규칙은 대개 드워킨이 계보 사항들이라고
부르는 유효성 기준들, 곧 법의 내용이 아니라 법들이 창설되거나 채택되는
방식과 형식에 관련되는 기준들을 제공함으로써 이 기능을 수행한다. 그러나
내가 말했듯이(250쪽), 승인 규칙은 그러한 계보 사항들에 더해 법들의
사실적 내용이 아니라 실체적 도덕 가치들(substantive moral values)이나
원칙들에 대한 그 부합과 관련되는 검사 기준들을 제공할 수도 있다. 물론
특정 사례들에서 판사들은 그러한 검사 기준들이 충족되는지 여부에 관해
의견불일치할 수 있고, 승인 규칙 안의 도덕적 검사 기준은 그러한
의견불일치를 해소하지 않을 것이다. 판사들은 특정 사례들에서
\textbf{(p.~259)} 그 검사 기준들이 무엇을 요구하는지에 관해서는
의견불일치하면서도, 그러한 검사 기준들의 관련성이 확립된 사법 실천에
의해 확정된 것이라는 데에는 동의할 수 있다. 이렇게 이해된 승인 규칙에는
규칙의 실천 이론이 온전히 적용 가능하다.

\subsubsection{\texorpdfstring{(ii) \emphksans{규칙들과 원칙들} (Rules and
Principles)}{(ii) 규칙들과 원칙들 (Rules and Principles)}}\label{ii-uxaddcuxce59uxb4e4uxacfc-uxc6d0uxce59uxb4e4-rules-and-principles}

오랫동안 이 책에 대한 드워킨의 비판 가운데 가장 잘 알려진 것은, 이 책이
법을 오직 ‘전부-아니면-무(all-or-nothing)’ 규칙들로만 이루어진 것으로
잘못 제시하고, 법적 추론과 재판에서 중요하고 독특한 역할을 하는 다른
종류의 법적 기준, 곧 법적 원칙들을 무시한다는 것이었다. 내 작업에서 이
결함을 발견한 일부 비판자들은 그것을 다소 고립된 흠으로 보았고, 내가
법체계의 구성 요소들로 법규칙들과 함께 법원칙들을 포함하기만 하면 그것을
고칠 수 있으며, 책의 주요 논제들 가운데 어느 것도 포기하거나 중대하게
수정하지 않고도 그렇게 할 수 있다고 생각했다. 그러나 이 비판 노선을 처음
강하게 제기한 드워킨은, 법원칙들은 나의 법이론 안에 포함될 수 있지만 그
대가는 그 중심 교설들의 포기라고 주장해 왔다. 그에 따르면 내가 법이
부분적으로 원칙들로 이루어져 있음을 인정한다면, 나는 한 체계의 당해 법이
법원들의 실천에서 수용된 승인 규칙이 제공하는 기준들에 의해 식별된다는
주장도, 현존하는 명시적 법이 결정을 지시하지 못하는 경우 법원들이
진정한, 비록 간극적인, 법창출 권한 또는 재량을 행사한다는 주장도, 법과
도덕 사이에 중요한 필연적 또는 개념적 연결이 없다는 주장도 일관되게
유지할 수 없다. 이 교설들은 내 법이론의 중심일 뿐 아니라 현대
법실증주의의 핵심을 구성하는 것으로 자주 받아들여지므로, 그것들의 포기는
상당한 중대성을 가진 문제일 것이다.

이 응답의 이 절에서 나는 내가 법원칙들을 무시했다는 비판의 여러 측면들을
검토하고, 그 비판에서 타당한 것이 무엇이든 내 이론 전체에 중대한 귀결
없이 수용될 수 있음을 보이고자 한다. 그러나 나는 이제, 재판과 법적
추론이라는 주제에 관해, 특히 내 비판자들이 법원칙들이라고 부르는
것들로부터의 논변들에 관해 이 책에서 내가 너무 적게 말했다는 점을 기꺼이
인정한다. 원칙들이 이 책에서 지나가듯 다루어졌을 뿐이라는 점은
결함이라는 데 나는 이제 동의한다.

그러나 정확히 내가 무시했다고 고발되는 것은 무엇인가? \textbf{(p.~260)}
법원칙들은 무엇이며, 그것들은 법규칙들과 어떻게 다른가? 법률가들이
사용하는 ‘원칙들’은 흔히 방대한 이론적·실천적 고려사항들을 포함하며, 그
가운데 일부만이 드워킨이 제기하려 한 쟁점들과 관련된다. ‘원칙’이라는
표현을, 사건들을 결정하는 법원들의 행위를 포함한 행위 기준들로 제한해
받아들인다 하더라도, 규칙들과 그러한 원칙들 사이의 대비를 긋는 방식들은
서로 다를 수 있다. 그러나 내가 원칙들을 무시했다고 비난한 모든
비판자들은 적어도 두 특징이 그것들을 규칙들과 구별한다는 데 동의할
것이라고 생각한다. 첫째는 정도의 문제이다. 원칙들은 규칙들에 비해 넓고
일반적이며 비특정적이다. 흔히 다수의 구별되는 규칙들로 여겨질 것들이
하나의 원칙의 예증들 또는 실례들로 제시될 수 있다는 의미에서 그렇다.
둘째 특징은, 원칙들이 어떤 목적, 목표, 권리자격(entitlement) 또는 가치를
어느 정도 명시적으로 참조하기 때문에, 어떤 관점에서는 유지하거나
고수하는 것이 바람직한 것으로 간주되며, 따라서 그것들을 예증하는
규칙들의 설명 또는 이론적 근거(rationale)를 제공할 뿐 아니라 적어도 그
정당화에도 기여하는 것으로 간주된다는 점이다.

규칙들과 관련하여 원칙들이 수행하는 설명적·정당화적 역할을 설명하는,
넓이와 어떤 관점에서의 바람직성이라는 이 두 비교적 논쟁의 여지가 없는
특징들 외에도, 세 번째 구별 특징이 있다. 나 자신은 이것이 정도의
문제라고 생각하지만, 드워킨은 그것을 핵심적이라고 본다. 그에 따르면
규칙들은 그것들을 적용하는 사람들의 추론에서 ‘전부-아니면-무 방식’으로
기능한다. 즉 어떤 규칙이 유효하고 주어진 사례에 적용 가능하다면, 그것은
법적 결과 또는 귀결을 ‘필연화한다’, 곧 확정적으로 결정한다.\footnote{{[}\emph{TRS}
  24면.{]}} 그가 법규칙들의 예로 든 것들 가운데에는 유료도로에서 시속
60마일의 최고속도를 정하는 규칙이나, 유언장의 작성·증명·효력을 규율하는
제정법들, 예컨대 두 증인이 서명하지 않으면 유언장은 유효하지 않다는
제정법 규칙이 있다. 드워킨에 따르면 법원칙들은 그러한 전부-아니면-무
규칙들과 다르다. 그것들은 적용될 때 결정을 ‘필연화하지’ 않고, 결정 쪽을
가리키거나 \textbf{(p.~261)} 결정에 유리하게 작용하거나, 압도될 수
있지만 법원들이 어느 한 방향으로 기울게 하는 것으로 고려하는 이유를
진술하기 때문이다. 짧게 말해, 나는 원칙들의 이러한 특징을 그
‘비확정적(non-conclusive)’ 성격이라고 부르겠다. 드워킨이 든 그러한
비확정적 원칙들의 예들 가운데에는 ‘법원들은 소비자와 공적 이익들이
공정하게 다루어지는지 보기 위해 {[}자동차{]} 구매 계약들을 면밀히
검토해야 한다’는 것처럼 상대적으로 특정적인 것들도 있고,\footnote{\emph{TRS}
  24면. 다음 판례에서 인용함: \emph{Henningsen v. Bloomfield Motors,
  Inc.}, 32 NJ 358, 161 A.2d 69 (1960), 387면, 85면.} ‘누구도 자신의
그른 행위로부터 이익을 얻어서는 안 된다’는 것처럼 훨씬 넓은 범위를
가지는 것들도 있다.\footnote{\emph{TRS} 25--26면.} 실제로 미국 헌법
제1·5·14 수정조항의 규정들처럼 미국 연방의회의 권한들과 주 입법에 대한
가장 중요한 헌법적 제한들의 상당수는 비확정적 원칙들로
기능한다.\footnote{{[}드워킨은 \emph{TRS} 27면에서 수정헌법 제1조가
  규칙(rule)인지 원칙(principle)인지 논의함.{]}} 드워킨에 따르면
법원칙들은 유효성의 차원이 아니라 \emphksans{무게}의 차원을 가진다는 점에서
규칙들과 다르다.\footnote{{[}\emph{TRS} 26면.{]}} 그래서 더 큰 무게의
다른 원칙과 충돌하면 한 원칙은 압도되어 결정을 확정하지 못할 수 있지만,
그럼에도 온전하게 존속하여 다른 사례들에서 더 작은 무게의 다른 원칙과
경쟁할 때 이길 수도 있다. 반면 규칙들은 유효하거나 유효하지 않을 뿐 이
무게의 차원을 가지지 않는다. 그래서 처음 정식화된 규칙들이 충돌하면,
드워킨에 따르면 그 가운데 오직 하나만 유효할 수 있고, 다른 규칙과의
경쟁에서 진 규칙은 그 경쟁 규칙과 일관되도록 다시 정식화되어야 하며
따라서 주어진 사례에는 적용될 수 없게 된다.\footnote{\emph{TRS}
  24--27면.}

나는 법원칙들과 법규칙들 사이의 이 날카로운 대비도, 유효한 규칙이 주어진
사례에 적용 가능하다면 원칙과 달리 항상 그 사례의 결과를 결정해야 한다는
견해도 받아들일 이유가 없다고 본다. 법체계가, 어떤 유효한 규칙은 그것이
적용 가능한 사례들에서 결과를 결정하지만, 더 중요하다고 판단되는 다른
규칙이 같은 사례에도 적용되는 경우에는 예외라고 승인하지 못할 이유는
없다. 따라서 주어진 사례에서 \textbf{(p.~262)} 더 중요한 규칙과의
경쟁에서 패배한 규칙은 원칙처럼, 다른 경쟁 규칙보다 더 중요하다고
판단되는 다른 사례들에서는 결과를 결정하기 위해 존속할 수
있다.\footnote{이 중요한 지점은 Raz와 Waluchow가 강조했지만, 나는 이에
  충분한 주의를 환기하지 못했다. J. Raz, “Legal Principles and the
  Limits of the Law”, \emph{Yale Law Journal} 제81권 (1972) 823면,
  832--834면; W. J. Waluchow, “Herculean Positivism”, \emph{Oxford
  Journal of Legal Studies} 제5권 (1985) 187면, 189--192면 참조.}

그러므로 드워킨에게 법은 전부-아니면-무 규칙들과 비확정적 원칙들 양쪽을
포함하며, 그는 이 차이를 정도의 문제로 보지 않는다. 그러나 나는 드워킨의
입장이 정합적일 수 있다고 생각하지 않는다. 그의 가장 이른 예들은
규칙들이 원칙들과 충돌할 수 있고, 어떤 원칙은 때로 규칙과의 경쟁에서
이기고 때로 질 것임을 함축한다. 그가 인용한 사례들에는 \emph{Riggs} v.
\emph{Palmer}가 포함된다.\footnote{115 N.Y. 506, 22 N.E. 188 (1889);
  \emph{TRS} 23면; \emph{LE} 15면 이하 참조.} 그 사건에서 사람은 자신의
그른 행위로부터 이익을 얻도록 허용되어서는 안 된다는 원칙은, 유언의
효력을 지배하는 제정법 규칙들의 명확한 문언에도 불구하고, 살인자가
희생자의 유언에 따라 상속하는 것을 배제하는 것으로 판시되었다. 이것은
원칙이 규칙과의 경쟁에서 이기는 예이지만, 그러한 경쟁의 존재 자체가
규칙들이 전부-아니면-무 성격을 가지지 않음을 분명히 보여 준다. 규칙들은
자신보다 더 무거울 수 있는 원칙들과 그러한 충돌 속으로 들어갈 수 있기
때문이다. 우리가 그러한 사례들을 (드워킨이 때때로 시사하듯이) 규칙들과
원칙들 사이의 충돌이 아니라, 문제된 규칙을 설명하고 정당화하는 원칙과
다른 어떤 원칙 사이의 충돌로 서술하더라도, 전부-아니면-무 규칙들과
비확정적 원칙들 사이의 날카로운 대비는 사라진다. 이 견해에서는 어떤
규칙이 그 문언에 따르면 적용 가능한 사례에서 결과를 확정하지 못하는 것은
그 정당화 원칙이 다른 원칙에 의해 무게에서 압도되기 때문이다. (드워킨이
또한 시사하듯이) 우리가 어떤 원칙을 명확히 정식화된 법규칙에 대한 새로운
해석을 위한 이유를 제공하는 것으로 생각하는 경우에도
마찬가지이다.\footnote{{[}드워킨의 논의는 \emph{TRS} 22--28면, \emph{LE}
  15--20면 참조.{]}}

법체계가 전부-아니면-무 규칙들과 비확정적 원칙들 양쪽으로 이루어진다는
주장 안의 이 비정합성은, 그 구별이 정도의 문제임을 인정하면 치유될 수
있다. \textbf{(p.~263)} 확실히, 적용 조건들의 충족이 몇몇 사례들, 곧 그
규정들이 더 중요하다고 판단되는 다른 규칙과 충돌할 수 있는 사례들을
제외하고는 법적 결과를 확정하기에 충분한 거의 확정적인 규칙들과, 단지
결정을 향해 가리키지만 매우 자주 그것을 확정하지 못할 수 있는 일반적으로
비확정적인 원칙들 사이에는 합리적 대비가 이루어질 수 있다.

나는 그러한 비확정적 원칙들로부터의 논변들이 재판과 법적 추론의 중요한
특징이며, 적절한 용어법으로 표시되어야 한다고 확실히 생각한다. 그것들의
중요성과 법적 추론에서의 역할을 보이고 예시한 점에서 드워킨에게는 큰
공이 있다. 그리고 내가 그것들의 비확정적 힘을 강조하지 않은 것은 분명
중대한 실수였다. 그러나 나는 ‘규칙’이라는 말을 사용하면서 법체계들이
오직 ‘전부-아니면-무’ 규칙들 또는 거의 확정적인 규칙들만 포함한다고
주장하려 의도하지 않았다. 나는 이 책 130--3쪽에서 내가 (아마도 썩
적절하지 않게) ‘가변적 법적 기준들’이라고 부른 것들, 곧 고려되어 다른
것들과 저울질되어야 할 요소들을 특정하는 기준들에 주의를 환기했을 뿐
아니라, 왜 어떤 행위 영역들은 ‘상당한 주의’와 같은 가변적 기준들이
아니라 드문 사례들을 제외하고 동일한 특정 행위들을 금지하거나 요구하는
거의 확정적인 규칙들에 의한 규율에 적합한지를 설명하려고도 했다(133--4쪽
참조). 그래서 우리는 살인과 절도에 반대하는 규칙들을 가지고 있는 것이지,
인간 생명과 재산에 대한 적절한 존중을 요구하는 원칙들만 가지고 있는 것이
아니다.

\subsection{4. 원칙들과 승인 규칙 (PRINCIPLES AND THE RULE OF
RECOGNITION)}\label{uxc6d0uxce59uxb4e4uxacfc-uxc2b9uxc778-uxaddcuxce59-principles-and-the-rule-of-recognition}

\subsubsection{\texorpdfstring{\emphksans{계보와 해석} (Pedigree and
Interpretation)}{계보와 해석 (Pedigree and Interpretation)}}\label{uxacc4uxbcf4uxc640-uxd574uxc11d-pedigree-and-interpretation}

드워킨은 법원들의 실천에서 드러나는 승인 규칙이 제공하는 기준들에 의해
법원칙들이 식별될 수 없으며, 원칙들은 법의 필수 요소들이므로 승인 규칙
교설은 폐기되어야 한다고 주장해 왔다. 그에 따르면 법원칙들은 오직 구성적
해석에 의해서만, 곧 한 법체계의 확립된 법 전체의 제도적 역사에 가장 잘
부합하고 그것을 가장 잘 정당화하는 유일한 원칙 집합의 구성원들로서만
식별될 수 있다. 물론 영국이든 미국이든 어떤 법원도 \textbf{(p.~264)}
당해 법을 식별하기 위해 그러한 체계-전체적 총체 기준을 명시적으로 채택한
적은 없으며, 드워킨은 자신의 신화적 이상 판사 ‘헤라클레스’와 구별되는
실제 인간 판사는 누구도 한 번에 자기 나라의 모든 법에 대한 해석을
구성하는 위업을 성취할 수 없다는 점을 인정한다. 그럼에도 그의 견해에서
법원들은 제한된 방식으로 ‘헤라클레스를 모방하려’ 하는 것으로 이해될 때
가장 잘 해명되며, 그는 자신들의 판결을 이렇게 바라보는 일이 ‘숨겨진
구조’를 드러내는 데 기여한다고 생각한다.\footnote{\emph{LE} 265면.}

제한된 형식의 구성적 해석에 의해 원칙들이 식별되는 가장 유명한 예, 영국
법률가들에게 익숙한 예는 \emph{Donoghue} v. \emph{Stevenson} 사건에서
애킨 경이 이전에는 정식화되지 않았던 ‘이웃 원칙(neighbour principle)’을,
여러 상황에서 주의책무를 확립하는 다양한 개별 규칙들의 밑바탕에 놓인
원칙으로 정식화한 것이다.\footnote{{[}1932{]} A.C. 562.}

나는 그러한 제한된 구성적 해석의 수행들에서 판사들이 헤라클레스의 총체적
체계-전체 접근을 모방하려 한다고 이해하는 것이 그럴듯하다고 보지 않는다.
그러나 여기서 나의 현재 비판은, 구성적 해석에 대한 몰두가 드워킨으로
하여금 많은 법원칙들이 그 지위를 확립된 법의 해석으로 기능하는 그 내용
덕분이 아니라, 그가 그 원칙들의 ‘계보’라고 부르는 것, 곧 인정된 권위
있는 원천에 의해 그것들이 창설되거나 채택된 방식 덕분에 얻는다는 사실을
무시하게 만들었다는 것이다. 이 몰두는 실제로 그를 이중의 오류로
이끌었다고 생각한다. 첫째는 법원칙들이 그 계보에 의해 식별될 수 없다는
믿음이고, 둘째는 승인 규칙이 오직 계보 기준들만 제공할 수 있다는
믿음이다. 이 두 믿음은 모두 잘못이다. 첫 번째 믿음이 잘못인 것은
원칙들의 비확정적 성격에도, 그 밖의 특징들에도, 그것들이 계보 기준들에
의해 식별되는 것을 배제하는 것은 전혀 없기 때문이다. 분명 성문헌법이나
헌법 수정조항 또는 제정법의 한 규정은 원칙들에 특징적인 비확정적
방식으로 작동하도록 의도된 것으로, 곧 어떤 다른 규칙이나 원칙이 대안적
결정을 위한 더 강한 이유들을 제시하는 사례들에서는 압도될 수 있는 결정
이유들을 제공하는 것으로 받아들여질 수 있다. 드워킨 자신도 의회가 표현의
자유를 제한해서는 안 된다고 규정하는 미국 헌법 제1수정조항이 바로 그런
방식으로 \textbf{(p.~265)} 해석되어야 한다고 상정했다.\footnote{{[}\emph{TRS}
  27면 참조.{]}} 또한 자기 자신의 그른 행위로부터 이익을 얻어서는 안
된다는 원칙 같은 커먼로의 몇몇 기본 원칙들을 포함한 일부 법원칙들은,
반대 방향을 가리키는 이유들에 의해 어떤 사례들에서는 압도될 수 있음에도
고려되어야 할 결정 이유들을 제공하는 것으로, 여러 다른 사례 범위에서
법원들이 일관되게 원용해 왔다는 점에서 ‘계보’ 검사 기준에 의해 법으로
식별된다. 계보 기준들에 의해 식별되는 법원칙들의 이러한 예들 앞에서,
법의 일부로 원칙들을 포함하는 것이 승인 규칙 교설의 포기를 함축한다는
일반 논변은 성공할 수 없다. 사실 내가 아래에서 보이듯이, 그것들의 포함은
그 교설과 양립 가능할 뿐 아니라 실제로 그 교설의 수용을 요구한다. 승인
규칙이 제공하는 계보 기준들에 의해 적어도 일부 법원칙들이 법으로
‘포착’되거나 식별될 수 있다는 점이 인정된다면, 그리고 분명 인정되어야
한다면, 드워킨의 비판은 더 겸손한 주장으로 축소되어야 한다. 곧 그렇게
포착될 수 없는 법원칙들이 많이 있는데, 그것들은 너무 많거나, 너무
일시적이거나, 변화나 수정에 너무 열려 있거나, 법체계의 제도적 역사와
실천들에 가장 잘 부합하고 그것들을 가장 잘 정당화하는 정합적 원칙 체계에
속한다는 검사 기준 이외의 어떤 검사 기준을 참조해서도 법원칙들로 식별될
수 있게 하는 특징을 가지지 않는다는 주장이다. 첫눈에 이 해석주의적 검사
기준은 승인 규칙이 제공하는 기준의 대안이 아니라, 몇몇 비판자들이
주장했듯이,\footnote{예: E. P. Soper, “Legal Theory and the Obligation
  of a Judge”, \emph{RDCJ} 3면, 16면; J. Coleman, “Negative and
  Positive Positivism”, \emph{RDCJ} 28면; D. Lyons, “Principles,
  Positivism and Legal Theory”, \emph{Yale Law Journal} 제87권 (1977)
  415면.} 원칙들을 그 계보가 아니라 그 내용에 의해 식별하는 그러한
기준의 복합적 ‘연성-실증주의적’ 형식일 뿐인 듯하다. 그러한 해석 기준을
포함하는 승인 규칙이, 위 251쪽 이하에서 논의한 이유들 때문에, 드워킨에
따르면 실증주의자가 원할 법 식별의 확실성 정도를 확보할 수 없다는 것은
사실이다. 그럼에도 해석적 검사 기준이 법-식별의 관행적 양식의 일부였음을
보이는 것은 여전히 \textbf{(p.~266)} 그것의 법적 지위에 대한 좋은 이론적
설명일 것이다. 그러므로 법의 일부로 원칙들을 인정하는 것과 승인 규칙
교설 사이에는 드워킨이 주장하는 그런 양립불가능성이 분명 없다.

앞의 두 단락의 논변은, 드워킨의 논지와 달리 원칙들을 법의 일부로
수용하는 것이 승인 규칙 교설과 양립 가능하다는 점을 보이기에 충분하다.
설령 드워킨의 해석적 검사 기준이 그가 주장하듯 원칙들을 식별하기에
유일하게 적절한 기준이라고 하더라도 그렇다. 그러나 사실 더 강한 결론이
보증된다. 곧 법원칙들이 그러한 기준에 의해 식별되려면 승인 규칙이
필요하다는 결론이다. 왜냐하면 드워킨의 해석적 검사 기준에 의해 드러날
어떤 법원칙을 식별하기 위한 출발점은, 그 원칙이 부합하고 정당화하는 데
도움을 주는 확립된 법의 어떤 특정 영역이기 때문이다. 그러므로 그 기준의
사용은 확립된 법의 식별을 전제하며, 그것이 가능하려면 법의 원천들과 그들
사이에 성립하는 우위와 종속의 관계들을 특정하는 승인 규칙이 필요하다.
\emph{Law's Empire}의 용어법으로 말하면, 밑바탕에 놓인 또는 암시적
법원칙들을 식별하는 해석 과제의 출발점들을 구성하는 법규칙들과 실천들은
‘전해석적 법(preinterpretive law)’을 구성하며, 드워킨이 그것에 관해
말하는 많은 것은, 그것의 식별을 위해 이 책에서 서술된 것처럼 법의 권위
있는 원천들을 식별하는 승인 규칙과 매우 유사한 것이 필요하다는 견해를
지지하는 듯하다. 여기서 내 견해와 드워킨의 견해 사이의 주된 차이는, 내가
법의 원천들을 식별하기 위한 기준들에 관해 판사들 사이에서 발견되는
일반적 합의를 그러한 기준들을 제공하는 \emphksans{규칙들}의 공유된 수용에
귀속시키는 반면, 드워킨은 규칙들이 아니라 같은 해석 공동체의 구성원들이
공유하는 ‘합의’,\footnote{{[}\emph{LE} 65--66면, 91--92면.{]}}
‘패러다임들’,\footnote{{[}\emph{LE} 72--73면.{]}} ‘가정들’\footnote{{[}\emph{LE}
  47, 67면.{]}}에 관해 말하기를 선호한다는 점이다. 물론 드워킨이 분명히
했듯이, 합의의 각 당사자가 일치하는 데 가지는 이유의 일부로 다른 이들의
일치가 들어 있지 않은 독립적 확신들의 합의와, 그것이 그러한 일부인
관행의 합의 사이에는 중요한 구별이 있다. 승인 규칙은 내 책에서 분명
\textbf{(p.~267)} 관행적 형식의 사법적 합의에 기초하는 것으로 취급된다.
그것이 그렇게 기초해 있다는 점은 적어도 영국법과 미국법에서는 아주
분명해 보인다. 영국 판사가 의회의 입법을, 또는 미국 판사가 헌법을, 다른
원천들에 대해 최고성을 갖는 법의 원천으로 취급하는 이유에는 분명 그의
사법 동료들이 그들의 선임자들이 그랬던 것처럼 이에 동의한다는 사실이
포함되기 때문이다. 기실 드워킨 자신도 입법 최고성 교설을 판사의 확신이
수행할 수 있는 역할을 제한하는 법적 역사의 날것의 사실로
말하고,\footnote{{[}\emph{LE} 401면.{]}} ‘해석적 태도는 같은 해석
공동체의 구성원들이 “무엇이 실천의 일부로 간주되는가”에 관해 적어도
대략 동일한 가정들을 공유하지 않으면 존속할 수 없다’고
진술한다.\footnote{\emph{LE} 67면.} 그러므로 나는 드워킨이 말하는
규칙들과 ‘가정들’, ‘합의’, ‘패러다임들’ 사이에 어떤 차이들이 남아 있든,
법의 원천들에 대한 사법적 식별에 관한 그의 설명은 실질적으로 나의 설명과
같다고 결론짓는다.

그러나 나의 견해와 드워킨의 견해 사이에는 큰 이론적 차이들이 남아 있다.
드워킨은 법원칙들을 위한 자신의 해석적 검사 기준을, 어떤 법체계들에서 그
법원들의 수용에 의존하여 존재와 권위를 가지는 관행적 승인 규칙이 취하는
특정 형식으로 다루는 나의 처리를 분명 거부할 것이다. 그의 견해에서는
이것이 법을 최선의 도덕적 빛 아래 보이도록 설계된 ‘구성적’ 해석의 기획,
곧 드워킨의 견해에서 법의 식별에 관여하는 기획을 완전히 잘못 제시하고
격하하는 것이기 때문이다. 그는 이러한 해석 양식을, 특정 법체계들의
판사들과 법률가들이 수용한 단순한 관행적 규칙이 요구하는 법-식별의
방법으로 구상하지 않는다. 대신 그는 그것을 법 이외에도 많은 사회사상과
사회적 실천의 중심 특징으로, 그리고 문학비평에서 이해되는 해석과 심지어
자연과학에서 이해되는 해석까지 포함하여 ‘모든 형식의 해석 사이의 깊은
연결’을 보여 주는 것으로 제시한다.\footnote{\emph{LE} 53면.} 그러나 이
해석 기준이 단지 관행적 규칙이 요구하는 법-식별의 양식이 아니고 다른
분과들에서 이해되는 해석과 친연성 및 연결을 가진다 하더라도, 드워킨의
총체적 해석 기준이 \textbf{(p.~268)} 실제로 법원칙들을 식별하는 데
사용되는 법체계들이 있다면 그 체계들에서 그 기준이 관행적 승인 규칙에
의해 제공될 수 있음은 여전히 사실이다. 그러나 이 완전한 총체 기준이
사용되는 실제 법체계들은 없고, \emph{Donoghue} v. \emph{Stevenson} 같은
사례들에서 잠재적 법원칙들을 식별하기 위해 더 겸손한 구성적 해석
수행들이 이루어지는 영국법과 미국법 같은 체계들만 있을 뿐이다. 그러므로
고려되어야 할 유일한 질문은 그러한 수행들이 관행적 승인 규칙이 제공하는
기준의 적용으로 이해되어야 하는지, 아니면 다른 어떤 방식으로 이해되어야
하는지, 그리고 그렇다면 그것들의 법적 지위가 무엇인지이다.

\subsection{5. 법과 도덕 (LAW AND
MORALITY)}\label{uxbc95uxacfc-uxb3c4uxb355-law-and-morality}

\subsubsection{\texorpdfstring{(i) \emphksans{권리들과 책무들} (Rights and
Duties)}{(i) 권리들과 책무들 (Rights and Duties)}}\label{i-uxad8cuxb9acuxb4e4uxacfc-uxcc45uxbb34uxb4e4-rights-and-duties}

이 책에서 나는, 법과 도덕 사이에는 서로 다른 많은 우연적 연결들이 있지만
법의 내용과 도덕 사이에는 필연적 개념적 연결들이 없으며, 따라서
도덕적으로 사악한 규정들도 법규칙들 또는 법원칙들로서 유효할 수 있다고
논변한다. 이러한 법과 도덕의 분리 형식의 한 측면은, 어떤 도덕적 정당화나
힘도 전혀 가지지 않는 법적 권리들과 책무들이 있을 수 있다는 점이다.
드워킨은 이 관념을 거부하고, 법적 권리들과 책무들의 존재에 관한
단언들에는 적어도 일응의 도덕적 근거들이 있어야 한다는 견해를 선호한다.
그래서 그는 ‘법적 권리들은 도덕적 권리들의 {[}한{]} 종으로 이해되어야
한다’는 관념을 자신의 법이론에서 ‘핵심적(crucial)’\footnote{\emph{RDCJ}
  260면.} 요소로 여기며, 그에 반대되는 실증주의 교설은 도덕적 근거 또는
힘이 전혀 없는 법적 권리들과 책무들이 있을 수 있다는 것을 분석 이전에
그저 아는 것으로 우리에게 주어지는 ‘법적 본질주의의 특이한
세계’\footnote{\emph{RDCJ} 259면.}에 속한다고 말한다. 일반 기술적
법리학이 법의 이해에 할 수 있는 기여의 종류를 이해하기 위해서는, 그의
일반 해석이론의 이점이 무엇이든, 법적 권리들과 책무들이 도덕적 힘이나
정당화를 결여할 수 있다는 교설에 대한 드워킨의 비판이 잘못되었음을 보는
것이 중요하다고 나는 생각한다. 그 이유는 다음과 같다. \textbf{(p.~269)}
법적 권리들과 책무들은 법이 그 강제적 자원들로 각각 개인의 자유를
보호하고 제한하거나, 개인들에게 법의 강제적 장치를 이용할 권한을
부여하거나 부인하는 지점이다. 그러므로 법들이 도덕적으로 좋든 나쁘든,
정의롭든 부정의하든, 권리들과 책무들은 법의 작용에서 인간들에게 최고로
중요한 초점들로서, 법들의 도덕적 이점들과 독립적으로 주목을 요구한다.
따라서 법적 권리들과 책무들의 진술들이 현실 세계에서 의미를 가질 수
있으려면 그 존재를 단언할 어떤 도덕적 근거가 있어야만 한다는 말은 참이
아니다.

\subsubsection{\texorpdfstring{(ii) \emphksans{법의 식별} (The Identification
of the
Law)}{(ii) 법의 식별 (The Identification of the Law)}}\label{ii-uxbc95uxc758-uxc2dduxbcc4-the-identification-of-the-law}

이 책에서 전개된 법이론과 드워킨의 이론 사이에서 법과 도덕의 연결들에
관련된 가장 근본적인 차이는 법의 식별에 관한 것이다. 내 이론에 따르면
법의 존재와 내용은 법의 사회적 원천들, 예컨대 입법, 사법 결정들, 사회
관습들을 참조하여 식별될 수 있으며, 그렇게 식별된 법이 그 자체로 법의
식별을 위한 도덕적 기준들을 편입한 경우를 제외하면 도덕에 대한 참조 없이
그렇게 될 수 있다. 반면 드워킨의 해석이론에서는 어떤 주제에 관한 당해
법이 무엇인지 진술하는 모든 법명제가 필연적으로 도덕 판단을 포함한다.
그의 총체적 해석이론에 따르면 법명제들은 다른 전제명제들과 함께, 법의
사회적 원천들을 참조하여 식별된 한 체계의 확립된 법 전체에 가장 잘
부합하고 그것에 대한 최선의 도덕적 정당화를 제공하는 원칙들의 집합에서
따라 나올 때에만 참이기 때문이다. 그러므로 이 전체적 총체 해석이론은
이중 기능을 가진다. 그것은 법을 식별하는 데에도, 그 법에 대한 도덕적
정당화를 제공하는 데에도 쓰인다.

이것이 \emph{Law's Empire}에서 ‘해석적’ 법과 ‘전해석적’ 법 사이의 구별을
도입하기 전 드워킨의 이론을 간략히 요약한 모습이었다. 법의 존재와 내용이
도덕에 대한 참조 없이 식별될 수 있다는 실증주의자의 이론에 대한 대안으로
고려할 때, 원래 상태의 드워킨 이론은 다음 비판에 취약했다. 사회적
원천들에 대한 참조로 식별된 법이 도덕적으로 사악한 곳에서는,
\textbf{(p.~270)} 그것을 가장 잘 ‘정당화’하는 원칙들은 그 법에 부합하는
원칙들 가운데 가장 덜 사악한 것들일 수밖에 없다. 그러나 그러한 가장 덜
사악한 원칙들은 어떤 정당화하는 힘도 가질 수 없으며, 무엇이 법으로
간주될 수 있는지에 어떤 도덕적 한계나 제약도 구성할 수 없다. 그리고
그것들이 아무리 사악한 법체계에도 부합할 수밖에 없기 때문에, 그것들을
참조하여 법을 식별한다고 주장하는 이론은 도덕에 대한 어떤 참조도 없이
법이 식별될 수 있다는 실증주의 이론과 구별되지 않는다. 드워킨이 ‘배경
도덕(background morality)’이라고 부른 것의 기준상 도덕적으로 건전한
원칙들, 단지 법에 부합하는 원칙들 가운데 도덕적으로 가장 건전한 것이
아닌 원칙들은 기실 무엇이 법으로 간주될 수 있는지에 도덕적 한계나 제약을
제공할 수 있다.\footnote{{[}\emph{TRS} 112, 128면, \emph{TRS} 93면
  참조.{]}} 나는 그 명제에 어떤 식으로도 반대하지 않지만, 그것은 법이
도덕에 대한 참조 없이 식별될 수 있다는 나의 주장과 완전히 양립 가능하다.

드워킨은 후기에 해석적 법과 전해석적 법 사이의 구별을 도입하면서, 우리가
도덕적으로 수용 가능하다고 볼 수 있는 그 법들의 어떤 해석도 불가능할
정도로 사악한 법체계들이 있을 수 있음을 인정한다. 그러한 경우에는 그가
설명하듯이 ‘내부 회의주의’\footnote{\emph{LE} 78--79면.}에 의지하여
그러한 체계들이 법임을 부정할 수 있다. 그러나 그러한 상황들을 서술하기
위한 우리의 자원들은 매우 유연하기 때문에, 우리는 대신 아무리 사악한
법체계들도 전해석적 의미에서는 법이라고 말할 수 있을 때 반드시 그 결론에
이르도록 구속되어 있지 않다.\footnote{{[}\emph{LE} 103면.{]}} 그러므로
우리는 가장 나쁜 나치 법들에 대해서조차 그것들이 법이 아니라고 말하도록
강제되지 않는다. 그 법들은 법창출의 형식들, 재판과 집행의 형식들과 같은
법의 많은 독특한 특징들을 도덕적으로 수용 가능한 체제들의 법들과
공유하면서, 오직 그 사악한 도덕적 내용에서만 그것들과 다를 수 있기
때문이다. 많은 맥락들과 많은 목적들에서는 그 도덕적 차이를 무시하고,
실증주의자와 함께 그러한 사악한 체계들이 법이라고 말할 충분한 이유들이
있을 수 있다. 이에 드워킨은 말하자면 자신의 해석적 관점에 대한 일반적
고수를 드러내는 단서, 곧 그러한 사악한 체계들은 오직 전해석적 의미에서만
법이라는 단서만을 덧붙일 것이다.

나는 우리 언어의 유연성에 대한 이 호소와 \textbf{(p.~271)} 이 지점에서
해석적 법과 전해석적 법 사이의 구별을 도입하는 일이 실증주의자의 논거를
약화하기보다 오히려 인정한다고 본다. 그것은, 그가 기술적 법리학에서는
법이 도덕에 대한 참조 없이 식별될 수 있다고 고집하지만, 확립된 법을
무엇이 가장 잘 정당화하는지에 관한 도덕 판단을 법의 식별이 항상
포함한다고 보는 정당화적 해석 법리학에서는 사정이 다르다는 메시지를
전달하는 것 이상의 일을 거의 하지 않기 때문이다. 물론 이 메시지는
실증주의자가 자신의 기술적 기획을 포기할 어떤 이유도 주지 않으며, 그렇게
하려는 의도도 없다. 그러나 이 메시지조차 한정되어야 한다. 법이 너무
사악하여 ‘내부 회의주의’가 적절한 경우에는 법의 해석이 어떤 도덕 판단도
포함하지 않으며, 드워킨이 이해하는 해석은 포기되어야 하기
때문이다.\footnote{{[}\emph{LE} 105면.{]}}

드워킨이 자신의 해석이론에 가한 또 하나의 수정은 그의 법적 권리 설명에
중요한 관련을 가진다. 원래 제시된 그의 총체 이론에서는 법의 식별과 그
정당화가 모두 한 체계의 확립된 법 모두에 가장 잘 부합하고 그것을 가장 잘
정당화하는 유일한 원칙 집합에서 따라 나오는 것으로 취급되었다. 그러므로
그러한 원칙들은, 내가 말했듯이, 이중 기능을 가진다. 그러나 한 체계의
확립된 법이 너무 사악해서 그 법에 대한 전반적 정당화 해석이 불가능할 수
있으므로, 드워킨은 이 두 기능이 분리될 수 있고, 도덕에 대한 어떤 참조도
없이 식별되는 법원칙들만 남을 수 있다고 관찰했다. 그러나 그러한 법은
드워킨이 모든 법적 권리들이 가진다고 주장하는 일응의 도덕적 힘을 가지는
어떤 권리들도 확립할 수 없다. 그럼에도 드워킨이 나중에 인정했듯이, 그
체계가 너무 사악해서 법 전체에 대한 어떤 도덕적 또는 정당화적 해석도
불가능한 곳에서도, 개인들이 적어도 일응의 도덕적 힘을 가진 권리들을
가진다고 적절히 말해질 수 있는 상황들이 여전히 있을 수 있다.\footnote{{[}\emph{LE}
  105--106면.{]}} 그것은 그 체계가, 예컨대 계약의 형성과 집행에 관련된
법들처럼 그 체계의 일반적 사악함에 영향을 받지 않을 수 있는 법들을
포함하고, 개인들이 자기 삶을 계획하거나 재산 처분들을 할 때 그러한
법들에 의존했을 수 있는 경우일 것이다. 그러한 상황들에 대응하기 위해
드워킨은, \textbf{(p.~272)} 일응의 도덕적 힘을 가진 법적 권리들과
책무들이 법에 대한 일반 해석이론에서 흘러나와야 한다는 자신의 원래
관념을 한정하고, 그러한 상황들이 자신의 일반 이론과 독립하여 개인들에게
어떤 도덕적 힘을 가진 법적 권리들을 귀속시키기 위한 ‘특별한 이유들’을
구성한다고 인정한다.

\subsection[6. 사법적 재량 (JUDICIAL DISCRETION)]{\texorpdfstring{6.
사법적 재량 (JUDICIAL
DISCRETION)\footnote{{[}본 절의 도입 문단의 대체 버전은 미주(endnote)에
  수록됨.{]}}}{6. 사법적 재량 (JUDICIAL DISCRETION)}}\label{uxc0acuxbc95uxc801-uxc7acuxb7c9-judicial-discretion73}

이 책의 법이론과 드워킨의 이론 사이의 가장 날카로운 직접 충돌은, 어떤
법체계에도 언제나 법적으로 규율되지 않은 일정한 사례들이 있을 것이며, 그
사례들에서는 어떤 지점에 관해 어느 쪽의 결정도 법에 의해 지시되지 않고
따라서 법이 부분적으로 불확정적이거나 불완전하다는 나의 논지로부터
생긴다. 그러한 사례들에서 판사가 결정을 내려야 하고, 벤담이 한때
주장했듯이 관할권을 부인하거나 현존 법에 의해 규율되지 않은 쟁점들을
입법부에 회부하여 결정하게 하지 않는다면, 판사는 이미 존재하는 확립된
법을 단지 적용하는 것이 아니라 자기 \emphksans{재량}을 행사하여 그 사례를
위한 법을 \emphksans{만들어야} 한다. 그러므로 그러한 법적으로 마련되지
않았거나 규율되지 않은 사례들에서 판사는 새로운 법을 만들면서, 동시에
자신의 법창출 권한들을 부여하고 제약하는 확립된 법을 적용한다.

법이 부분적으로 불확정적이거나 불완전하고 판사가 제한된 법창출 재량을
행사하여 그 틈들을 메운다는 이 구상은, 드워킨에 의해 법과 사법 추론
양쪽에 대한 오도적 설명으로 거부된다. 그는 사실상 불완전한 것은 법이
아니라 실증주의자의 법 그림이며, 이 점은 사회적 원천들을 참조하여
식별되는 \emphksans{명시적} 확립된 법 외에도 그 명시적 법에 가장 잘
부합하거나 정합하고 또한 그것에 대한 최선의 도덕적 정당화를 제공하는
원칙들인 \emphksans{암시적} 법원칙들을 포함하는 법에 대한 자신의 ‘해석적’
설명에서 드러난다고 주장한다. 이 해석적 견해에서는 법이 결코
불완전하거나 불확정적이지 않으므로, 판사는 결정을 내리기 위해 법
바깥으로 나가 법창출 권한을 행사할 계기를 결코 가지지 않는다. 따라서
\textbf{(p.~273)} 법의 사회적 원천들이 어떤 법적 쟁점에 관해 결정을
확정하지 못하는 ‘어려운 사례들’에서 법원들이 향해야 할 것은, 그 도덕적
차원들을 가진 그러한 암시적 원칙들이다.

법에 의해 부분적으로 규율되지 않은 채 남은 사례들을 규율하기 위해 내가
판사들에게 귀속시키는 법창출 권한들은 입법부의 권한들과 다르다는 점이
중요하다. 판사의 권한들은 입법부가 상당히 자유로울 수 있는 많은 제약들,
곧 \emphksans{그의 선택을 좁히는} 제약들에 종속될 뿐 아니라, 특정한 당해
사례들을 처리하기 위해서만 행사되므로 대규모 개혁이나 새로운 법전들을
도입하는 데 사용될 수 없다. 그래서 그의 권한들은 많은 실체적
제약들(substantive constraints)에 종속될 뿐 아니라 \emphksans{간극적}이다.
그럼에도 현존 법이 어떤 결정도 올바른 것으로 지시하지 못하는 지점들이
있을 것이며, 그런 경우들을 결정하기 위해 판사는 자신의 법창출 권한들을
행사해야 한다. 그러나 그는 이것을 자의적으로 해서는 안 된다. 곧 그는
언제나 자기 결정을 정당화하는 어떤 일반적 이유들을 가져야 하고, 자신의
믿음들과 가치들에 따라 결정함으로써 양심적 입법자가 하듯이 행위해야
한다. 그러나 그가 이 조건들을 충족한다면, 그는 법에 의해 지시되지 않고
유사한 어려운 사례들에 직면한 다른 판사들이 따른 것들과 다를 수 있는
결정 기준들이나 이유들을 따를 자격이 있다.

법에 의해 불완전하게 규율된 채 남은 사례들을 처리하기 위해 법원들이
그러한 제한된 재량적 권한을 행사한다는 나의 설명에 맞서, 드워킨은 세
가지 주요 비판을 제기한다. 첫째는 이 설명이 사법 절차와 ‘어려운
사례들’에서 법원들이 하는 일에 대한 거짓 기술이라는 것이다.\footnote{{[}\emph{TRS}
  81면; \emph{LE} 37--39면 비교.{]}} 이를 보이기 위해 드워킨은 판사들과
법률가들이 판사의 과제를 서술할 때 사용하는 언어와 사법적 의사결정의
현상학에 호소한다. 사건들을 결정하는 판사들과 자신들에게 유리하게
결정하라고 그들을 촉구하는 법률가들은, 새로운 사례들에서조차 판사가 법을
‘만든다’고 말하지 않는다는 것이다. 그러한 사례들 가운데 가장 어려운
것들에서조차 판사는, 실증주의자가 시사하듯이, 결정 과정에 완전히 다른 두
단계가 있다는 자각을 드러내지 않는 경우가 많다. 곧 판사가 먼저 현존 법이
어느 쪽의 결정도 지시하지 못한다고 발견하는 단계와, 그 다음 현존 법에서
돌아서서 무엇이 최선인지에 관한 자신의 관념에 따라 당사자들을 위해
\emphksans{처음부터} 그리고 \emphksans{사후적으로} 법을 만드는 단계 말이다.
\textbf{(p.~274)} 대신 법률가들은 판사가 언제나 현존 법을 발견하고
집행하는 데 관여하는 것처럼 판사에게 말하고, 판사는 법이 모든 사례에
대한 해법이 판사의 발명이 아니라 발견을 기다리는, 빈틈없는 권리자격들의
체계인 것처럼 말한다.

사법 절차의 익숙한 수사가 발전된 법체계에는 법적으로 규율되지 않은
사례들이 없다는 관념을 북돋는다는 데에는 의심의 여지가 없다. 그러나
이것은 얼마나 진지하게 받아들여져야 하는가? 물론 입법자와 판사의 구별을
극화하고, 판사는 현존 법이 명확할 때 그러하듯 언제나 자신이 만들거나
빚어내지 않는 법의 단순한 ‘입’이라고 주장하는 오랜 유럽 전통과 권력분립
교설이 있다. 그러나 법정에서 사건들을 결정할 때 판사들과 법률가들이
사용하는 의례적 언어와, 사법 절차에 관한 그들의 더 성찰적인 일반
진술들을 구별하는 것이 중요하다. 미국의 올리버 웬델 홈스와 카도조,
영국의 맥밀런 경이나 래드클리프 경 또는 리드 경과 같은 격의 판사들,
그리고 학계와 실무계 양쪽의 수많은 다른 법률가들은, 법에 의해 불완전하게
규율된 채 남은 사례들이 있어서 판사에게 피할 수 없는, 비록 ‘간극적’인,
법창출 과제가 있으며, 법에 관한 한 많은 사례들이 어느 쪽으로든 결정될 수
있다고 주장해 왔다.

하나의 주요한 고려사항은 판사들이 때로 법을 만들면서 동시에 적용한다는
주장에 대한 저항을 설명하는 데 도움을 주며, 사법적 법창출을 입법부의
법창출과 구별하는 주요 특징들도 해명한다. 그것은 규율되지 않은 사례들을
결정할 때 법원들이, 자신들이 만드는 새로운 법이 비록 \emphksans{새로운}
법이기는 하지만 이미 현존 법에 발판을 가진 것으로 인정된 원칙들이나
밑받침 이유들과 부합하도록 보장하기 위해 유비에 따라 진행하는 데
전형적으로 부여하는 중요성이다. 특정 제정법들이나 선례들이 불확정적인
것으로 드러나거나 명시적 법이 침묵할 때, 판사들이 단지 법서들을 밀쳐
두고 법으로부터 더 이상의 지침 없이 입법하기 시작하지 않는 것은
사실이다. 그러한 사례들을 결정하면서 그들은 매우 자주 현존 법의 상당한
관련 영역이 예증하거나 증진하는 것으로 이해될 수 있고 당해 어려운 사례에
대한 확정적 답을 향해 가리키는 어떤 일반 원칙이나 일반적 목표 또는
목적을 인용한다. \textbf{(p.~275)} 이것은 기실 드워킨 재판 이론의 매우
두드러진 특징인 ‘구성적 해석’의 바로 그 핵심이다. 그러나 이 절차는
사법적 법창출의 순간을 분명 유예하지만, 제거하지는 않는다. 어떤 어려운
사례에서도 경쟁하는 유비들을 뒷받침하는 서로 다른 원칙들이 나타날 수
있고, 판사는 법에 의해 자신에게 이미 규정된 우선성의 질서가 아니라
양심적 입법자처럼 무엇이 최선인지에 관한 자신의 감각에 의존하여 그들
사이에서 선택해야 하는 경우가 많기 때문이다. 그러한 모든 사례들에 대해,
경쟁하는 저차-질서 원칙들에 상대적 무게나 우선성을 배정하는 고차-질서
원칙들의 어떤 유일한 집합이 현존 법 안에서 언제나 발견될 수 있을 때에만,
사법적 법창출의 순간은 단순히 유예되는 데 그치지 않고 제거될 것이다.

드워킨의 나의 사법 재량 설명에 대한 다른 비판들은 그것을 기술적으로
거짓이라고 비난하는 것이 아니라, 비민주적이고 부정의한 법창출 형식을
지지한다고 비난한다.\footnote{{[}\emph{TRS} 84--85면.{]}} 판사들은 보통
선출되지 않으며, 민주주의에서는 오직 선출된 국민의 대표자들만 법창출
권한들을 가져야 한다고 주장된다. 이 비판에는 많은 답변들이 있다. 법이
규율하지 못한 분쟁들을 처리하기 위해 판사들에게 법창출 권한들이 맡겨져야
한다는 것은, 입법부에 회부하는 것 같은 대안적 규율 방법들의 불편을
피하기 위해 치러야 할 필수적 대가로 여겨질 수 있다. 그리고 판사들이 이
권한들의 행사에서 제약되어 법전들이나 광범위한 개혁들을 형성할 수 없고
특정 사례들이 던지는 특정 쟁점들을 처리하기 위한 규칙들만 만들 수
있다면, 그 대가는 작아 보일 수 있다. 둘째, 제한된 입법 권한들을 행정부에
위임하는 것은 현대 민주주의들의 익숙한 특징이며, 그러한 위임이 사법부에
이루어진다고 해서 민주주의에 더 큰 위협으로 보이지는 않는다. 두 위임
형식 모두에서 선출된 입법부는 보통 잔여적 통제권을 가지며, 받아들일 수
없다고 보는 어떤 종속된 법들도 폐지하거나 수정할 수 있다. 미국에서처럼
입법부의 권한들이 성문헌법에 의해 제한되고 법원들이 광범위한 심사
권한들을 가지는 경우, \emphksans{민주적으로 선출된} 입법부가 어떤 사법적 입법
하나를 뒤집을 수 없음을 알게 될 수 있다는 것은 사실이다. 그러면
\textbf{(p.~276)} 궁극적 민주적 통제는 헌법 수정이라는 번거로운 장치를
통해서만 확보될 수 있다. 그것이 정부에 대한 법적 제약들을 위해 치러야
하는 대가이다.

드워킨은 사법적 법창출이 부정의하며, 물론 보통 부정의한 것으로 여겨지는
소급적 또는 \emphksans{사후적} 법창출의 한 형식으로 그것을 비난한다는 추가
고발을 한다. 그러나 소급적 법창출을 부정의한 것으로 보는 이유는, 행위할
때 자기 행위의 법적 귀결들이 그 행위 당시 확립되어 있던 알려진 법의
상태에 의해 결정될 것이라는 가정에 의존했던 사람들의 정당한 기대들을
그것이 좌절시키기 때문이다. 하지만 이 반론은, 분명히 확립된 법에 대한
법원의 소급적 변경이나 번복에 대해서는 힘을 가질 수 있다 하더라도,
어려운 사례들에서는 전혀 관련 없어 보인다. 그것들은 법이 불완전하게
규율한 채 남겨 둔 사례들이고, 기대들을 정당화할 명확히 확립된 알려진
법의 상태가 없는 경우들이기 때문이다.

\newpage

\subsection{후기 주석}\label{uxd6c4uxae30-uxc8fcuxc11d}


{\footnotesize
\noteentry 272쪽. {[}이 절의 대체 시작문이 여기에 포함되어 있으며, 폐기되지
않음.{]}

\begin{quote}
드워킨(Dworkin)은 오랜 시간에 걸친 재판(adjudication)에 관한 저술
전반에서, 법원들이 현존 법에 의해 불완전하게 규율된 채 남은 사례들을
결정할 법창출 권한(law-creating power)이라는 의미의 재량(discretion)을
가진다는 것을 부인하는 입장을 한결같이 견지해 왔다. 기실 그는, 사소한
예외를 제외하면 그러한 사례들은 존재하지 않으며, 유명하게 말했듯이 어떤
사례에서 생기는 어떤 법적 쟁점에 관해 당해 법이 무엇인지에 대한 의미
있는 질문에는 언제나 단 하나의 ‘정답(right answer)’이 있다고 논변해
왔다. \footnote{{[}그의 글 `No Right Answer?' in P. M. S. Hacker and J.
  Raz (eds.), \emph{Law, Morality and Society} (1977), pp.~58--84. 개정
  후 재수록: `Is There Really No Right Answer in Hard Cases?'
  \emph{AMP}, 제5장 참조.{]}}
\end{quote}

\begin{quote}
그러나 이러한 불변하는 듯한 교설의 외양에도 불구하고, 드워킨이 후기에
해석적 관념들(interpretive ideas)을 자신의 법이론에 도입하고 모든
법명제들(propositions of law)이 그가 이 표현에 부여한 특수한 의미에서
‘해석적’이라고 주장한 것은, 라즈가 처음 분명히 했듯이,\footnote{{[}J.
  Raz, `Dworkin: A New Link in the Chain', \emph{74 California Law
  Review}, 1103 (1986), pp.~1110, 1115--16 참조.{]}} 법원들이 실제로
법창출 재량을 가지며 자주 행사한다는 점을 인정하게 함으로써 이 입장의
실질을 나의 입장에 매우 가까이 가져왔다. 논변해 볼 수 있듯이, 해석적
관념들이 그의 이론에 도입되기 전에는, 드워킨의 강한 의미의 사법 재량
부정과 항상 정답이 있다는 그의 주장이 사건들을 결정하는 판사의 역할은
\emphksans{현존 법을 분별해 내고} \emphksans{집행하는 것}이라는 관념과 결합되어
있었기 때문에, 우리의 각 재판 설명 사이에는 큰 차이가 있는 듯했다.
그러나 물론 법원들이 사건들을 결정하면서 법창출 재량을 자주 행사한다는
나의 주장과 매우 날카롭게 충돌했던 이 초기 구상은 더 이상 이 절에서 전혀
등장하지 않는다.
\end{quote}

{[}섹션 6의 대체 시작문의 텍스트는 여기에서 끝난다.{]}

\subsection{후기 제3판
주석}\label{uxd6c4uxae30-uxc81c3uxd310-uxc8fcuxc11d}

이것은 아마도 어떤 후기(postscript)에 관해 책 한 권 전체가 나온 유일한
사례일 것이다:

Jules Coleman 엮음, \emph{Hart's Postscript}. 수록된 에세이들은
후기(postscript)뿐만 아니라 하트의 법이론 전반을 이해하는 데 매우
유익하다. 또한, Ronald Dworkin, “Hart's Postscript and the Character of
Political Philosophy” (\emph{Oxford Journal of Legal Studies} 제24권,
2004년, 1면)을 참조하라.
}
